%% 3-5 グラフの傾きは「変化の速さ」（第12回：組 V3-3。ページ制限なし。偶数ページ）  % noruby
%% 第5巻の思い出そう（5-0・5-1・5-3）が指す形：  % noruby
%%   「グラフの傾き＝たての変化量÷横の変化量＝Δy/Δx。右上がりなら正、右下がりなら負」  % noruby
%%   「傾きが m で、点 (x1, y1) を通る直線の式は y−y1＝m(x−x1)」（傾きを m と書くことも多い。y＝ax＋b の a と同じ）  % noruby
\makeatletter
% 式を、まん中に1行で置く  \bkthrfline{…}  % noruby
\newcommand\bkthrfline[1]{\par\addvspace{3pt}{\centering #1\par}\addvspace{3pt}\noindent\ignorespaces}
% 説明の段の見出し（丸数字＋太字の1行）  \bkthrfstep{番号}{見出し}  % noruby
\newcommand\bkthrfstep[2]{\par\addvspace{6pt}\Needspace{4\baselineskip}\noindent
  {\bkcircled[bkOrange]{#1}\enspace\sffamily\bfseries\color{bkOrange!75!black}#2}\par\nobreak\vspace{2pt}}
% 坂（直角三角形）  \bkthrfslope{(左下)}{横 mm}{高さ mm}{色}  % noruby
\newcommand\bkthrfslope[4]{\begin{scope}[shift={#1}]
  \fill[#4!22] (0,0) -- (#2,0) -- (#2,#3) -- cycle;
  \draw[bkGray, line width=0.5pt] (0,0) -- (#2,0) -- (#2,#3);
  \draw[#4!85!black, line width=1.2pt] (0,0) -- (#2,#3);
\end{scope}}
% 水そう（左下を原点に、はば22mm・高さ26mm。水の高さ #2 mm）  \bkthrftank{(x,y)}{水の高さ}  % noruby
\newcommand\bkthrftank[2]{\begin{scope}[shift={#1}]
  \fill[bkBlue!25] (0.4,0.4) rectangle (21.6,#2);
  \draw[bkBlue!80!black, line width=0.5pt] (0.4,#2) -- (21.6,#2);
  \draw[bkInk!80, line width=0.9pt] (0,26) -- (0,0) -- (22,0) -- (22,26);
  % じゃ口
  \fill[bkGray!70] (-6,30) rectangle (6,32.5);
  \fill[bkGray!70] (4,27.5) rectangle (6,32.5);
  \fill[bkBlue!45] (4.4,27.5) rectangle (5.6,{#2});
\end{scope}}
\makeatother

\bkchapter{3}{5}{グラフの\R{傾}{かたむ}きは「\R{変化}{へん|か}の\R{速}{はや}さ」}{\bklevel{2}}{グラフの\R{傾}{かたむ}きは「\R{変化}{へん|か}の\R{速}{はや}さ」}
\bkmichishirube{2-3,2-15,2-16,3-2,3-4}{3-6,3-7,5-1,5-3,8-2}

\begin{goal}
\begin{itemize}[label={\color{\bkTT}●}, leftmargin=1.6\zw, itemsep=0pt, topsep=1pt]
\item \bklevel{2}\enspace グラフの\R{傾}{かたむ}きを「たての\R{変化量}{へん|か|りょう}÷\R{横}{よこ}の\R{変化量}{へん|か|りょう}」で\R{求}{もと}め、\R{正}{せい}・\R{負}{ふ}・0の\R{意味}{い|み}がわかる。
\item \bklevel{2}\enspace 1\R{次関数}{じ|かん|すう} $y=ax+b$ の\R{傾}{かたむ}きと\R{切片}{せっ|ぺん}を\R{読}{よ}み、$x$-$t$ グラフの\R{傾}{かたむ}きを\R{速}{はや}さとして\R{使}{つか}える。
\item \bklevel{3}\enspace 1\R{点}{てん}と\R{傾}{かたむ}きから、\R{直線}{ちょく|せん}の\R{式}{しき} $y-y_1=m(x-x_1)$ を\R{作}{つく}れる。
\end{itemize}
\end{goal}

\begin{hatena}
\R{学校}{がっ|こう}の\R{近}{ちか}くに、2つの\R{坂}{さか}があります。\R{坂}{さか}Aは、\R{横}{よこ}に\qty{20}{m}\R{進}{すす}むあいだに\qty{2}{m}\R{上}{あ}がります。\R{坂}{さか}Bは、\R{横}{よこ}に\qty{5}{m}\R{進}{すす}むあいだに\qty{1}{m}\R{上}{あ}がります。\R{自転車}{じ|てん|しゃ}でのぼるとき、きついのはどっちかな？

\R{坂}{さか}の「\R{急}{きゅう}さ」を、1つの\R{数}{かず}で\R{言}{い}うには、どうすればよいでしょう。
\begin{bkfigure}
\begin{tikzpicture}[x=1mm, y=1mm]
  % 1 m を 3.6 mm でかく（横もたても同じ縮尺）
  \bkthrfslope{(0,0)}{72}{7.2}{bkBlue}
  \node[font=\sffamily\bfseries\small, text=bkBlue!80!black, anchor=south east] at (34,4.6) {\R{坂}{さか}A};
  \node[bk tick, anchor=north] at (36,-1.5) {\R{横}{よこ}\qty{20}{m}};
  \node[bk tick, anchor=west] at (73.5,3.6) {\qty{2}{m}};
  \bkthrfslope{(90,0)}{18}{3.6}{bkOrange}
  \node[font=\sffamily\bfseries\small, text=bkOrange!80!black, anchor=south east] at (98,3.2) {\R{坂}{さか}B};
  \node[bk tick, anchor=north] at (99,-1.5) {\R{横}{よこ}\qty{5}{m}};
  \node[bk tick, anchor=west] at (109.5,1.8) {\qty{1}{m}};
\end{tikzpicture}
\bkfigcaption{\R{横}{よこ}もたても、\qty{1}{m}を\R{同}{おな}じ\R{長}{なが}さでかいた\R{図}{ず}。}
\end{bkfigure}
\end{hatena}

\noindent グラフの\R{線}{せん}も、\R{坂}{さか}のように\R{上}{あ}がったり\R{下}{さ}がったりします。3-4でかいた $x$-$t$ グラフでは、\R{線}{せん}が\R{急}{きゅう}な\R{所}{ところ}ほど、\R{位置}{い|ち}が\R{速}{はや}く\R{変}{か}わっていました。この\R{章}{しょう}では、\R{坂}{さか}やグラフの「\R{急}{きゅう}さ」を\R{数}{かず}にします。その\R{数}{かず}が、「\R{変化}{へん|か}の\R{速}{はや}さ」を\R{表}{あらわ}していることを\R{確}{たし}かめます。

%% ---------------- やってみよう・思い出そう ----------------
\begin{yattemiyou}[スロープ・\R{階段}{かい|だん}・すべり\R{台}{だい}の\R{急}{きゅう}さをくらべよう]
\noindent\begin{minipage}[c]{0.58\linewidth}
\begin{enumerate}[label={\bkcircled[bkGreen]{\arabic*}}, leftmargin=2.4\zw, itemsep=1pt, topsep=0pt]
\item メジャーで、\R{学校}{がっ|こう}の\R{入}{い}り\R{口}{ぐち}のスロープの「\R{横}{よこ}の\R{長}{なが}さ」と「\R{上}{あ}がる\R{高}{たか}さ」をはかろう。
\item \R{階段}{かい|だん}は1\R{段}{だん}の「おくゆき」と「\R{高}{たか}さ」を、すべり\R{台}{だい}はすべる\R{所}{ところ}の「\R{横}{よこ}の\R{長}{なが}さ」と「\R{高}{たか}さ」をはかろう。
\item \R{高}{たか}さ÷\R{横}{よこ}の\R{長}{なが}さを\R{電卓}{でん|たく}で\R{計算}{けい|さん}しよう。いちばん\R{急}{きゅう}なのはどれかな？
\end{enumerate}
\end{minipage}\hfill
\begin{minipage}[c]{0.38\linewidth}\centering
\begin{tikzpicture}[x=1mm, y=1mm]
  % 階段3段（おくゆき 30 cm・高さ 16 cm を、おくゆき 14 mm・高さ 7.5 mm でかく）
  \fill[bkGray!30] (0,0) -- (0,7.5) -- (14,7.5) -- (14,15) -- (28,15) -- (28,22.5) -- (42,22.5) -- (42,0) -- cycle;
  \draw[bkInk!80, line width=0.7pt] (-3,0) -- (0,0) -- (0,7.5) -- (14,7.5) -- (14,15) -- (28,15) -- (28,22.5) -- (42,22.5);
  \draw[bkRed, line width=1.2pt] (0,7.5) -- (14,7.5);
  \node[font=\sffamily\scriptsize, text=bkRed!80!black, anchor=south, inner sep=1pt] at (6.5,8.4) {おくゆき};
  \draw[bkBlue, line width=1.2pt] (28,15) -- (28,22.5);
  \node[font=\sffamily\scriptsize, text=bkBlue!80!black, anchor=east, inner sep=1pt] at (26.4,18.8) {\R{高}{たか}さ};
  % メジャー
  \filldraw[fill=bkYellow!70, draw=bkYellow!60!black, line width=0.5pt, rounded corners=1mm] (2,-12) rectangle (12,-4);
  \fill[bkYellowBg] (12,-9.2) rectangle (42,-6.8);
  \draw[bkYellow!60!black, line width=0.4pt] (12,-9.2) rectangle (42,-6.8);
  \foreach \k in {0,...,14} \draw[bkInk, line width=0.3pt] (13+\k*2,-6.8) -- (13+\k*2,-7.8);
\end{tikzpicture}
\end{minipage}
\par\vspace{4pt}
ソラさんがはかった\R{結果}{けっ|か}です。
\begin{bkcenter}
\begin{bktable}{colspec={Q[l,m,wd=30mm] Q[c,m,wd=27mm] Q[c,m,wd=27mm] Q[c,m,wd=27mm]}, rowsep=1.5pt, row{1}={font=\sffamily\small}}
 & \R{横}{よこ}の\R{長}{なが}さ & \R{上}{あ}がる\R{高}{たか}さ & \R{高}{たか}さ÷\R{横}{よこ} \\
スロープ & \qty{360}{cm} & \qty{30}{cm} & \\
\R{階段}{かい|だん}1\R{段}{だん} & \qty{30}{cm} & \qty{16}{cm} & \\
すべり\R{台}{だい} & \qty{240}{cm} & \qty{120}{cm} & \\
\end{bktable}
\end{bkcenter}
\end{yattemiyou}

\begin{omoidasou}
\bkrecall{2-3}{\R{変化量}{へん|か|りょう}は「あと－まえ」。$\Delta x$ は「$x$ の\R{変化量}{へん|か|りょう}」。\R{速}{はや}さ＝$\Delta x\div\Delta t$。}
\bkrecall{2-15}{\R{比例}{ひ|れい} $y=ax$ のグラフは\R{原点}{げん|てん}を\R{通}{とお}る\R{直線}{ちょく|せん}。$a$ は1あたりの\R{量}{りょう}で、\R{大}{おお}きいほど\R{線}{せん}は\R{急}{きゅう}。}
\bkrecall{2-14}{\R{等式}{とう|しき}の\R{両辺}{りょう|へん}に\R{同}{おな}じ\R{数}{かず}をたしても、ひいても、かけても、\R{等式}{とう|しき}は\R{成}{な}り\R{立}{た}つ。}
\bkrecall{3-2}{\R{点}{てん} $(x,\ y)$ は、\R{原点}{げん|てん}から\R{右}{みぎ}へ $x$、\R{上}{うえ}へ $y$ の\R{所}{ところ}。$x_1$ は「エックス・いち」と\R{読}{よ}み、かけ\R{算}{ざん}ではない。}
\bkrecall{3-4}{\R{横}{よこ}\R{軸}{じく}を\R{時間}{じ|かん} $t$、たて\R{軸}{じく}を\R{位置}{い|ち} $x$ にしたグラフを、$x$-$t$ グラフという。}
\end{omoidasou}

\noindent やってみようの\R{答}{こた}え：スロープは $30\div360=0.083\cdots$ です（$\cdots$ は\R{数字}{すう|じ}がずっと\R{続}{つづ}くしるし）。\R{階段}{かい|だん}は $16\div30=0.533\cdots$、すべり\R{台}{だい}は $120\div240=0.5$ です。\R{数}{かず}が\R{大}{おお}きいほど\R{急}{きゅう}なので、いちばん\R{急}{きゅう}なのは\R{階段}{かい|だん}、いちばんゆるやかなのはスロープです。\R{階段}{かい|だん}は\R{急}{きゅう}でも、\R{段}{だん}があるので\R{足}{あし}がすべりません。

\bkfukidashi{sora}{\R{高}{たか}さだけを\R{見}{み}たら、すべり\R{台}{だい}の\qty{120}{cm}がいちばん\R{大}{おお}きいのに、\R{階段}{かい|だん}のほうが\R{急}{きゅう}なんだね。}

%% ---------------- 数学の道具 ----------------
\bksection{doguu}{\R{数学}{すう|がく}の\R{道具}{どう|ぐ}}

\bksubsection{1}{\R{急}{きゅう}さを\R{数}{かず}で\R{言}{い}う：\R{傾}{かたむ}き}

\leavevmode\label{bk:3-5:katamuki}\bkhatenaref の2つの\R{坂}{さか}をくらべます。まず、よくある\R{考}{かんが}え\R{方}{かた}をためしてみましょう。

\bkng{\R{坂}{さか}Aは\qty{2}{m}、\R{坂}{さか}Bは\qty{1}{m}\R{上}{あ}がる。だから\R{坂}{さか}Aのほうが\R{急}{きゅう}}
\noindent この\R{考}{かんが}え\R{方}{かた}では、\R{横}{よこ}に\R{進}{すす}む\R{長}{なが}さを\R{見}{み}ていません。\R{坂}{さか}Aは\qty{2}{m}\R{上}{あ}がるのに、\qty{20}{m}も\R{進}{すす}みます。\R{坂}{さか}Bは\qty{1}{m}\R{上}{あ}がるのに、たった\qty{5}{m}です。もし\R{高}{たか}さだけでくらべてよいなら、\R{長}{なが}い\R{長}{なが}い\R{道}{みち}を\qty{2}{m}だけ\R{上}{あ}がる、ほとんど\R{平}{たい}らな\R{道}{みち}も「\R{急}{きゅう}」になってしまいます。\R{急}{きゅう}さをくらべるには、\R{横}{よこ}の\R{長}{なが}さをそろえる\R{必要}{ひつ|よう}があります。

\bkthrfstep{1}{\R{方法}{ほう|ほう}1：\R{横}{よこ}の\R{長}{なが}さをそろえる}
\R{坂}{さか}Bを4つつなぐと、\R{横}{よこ}は $5\times4=20$ で\qty{20}{m}、\R{高}{たか}さは $1\times4=4$ で\qty{4}{m}になります。\R{同}{おな}じ\R{急}{きゅう}さの\R{坂}{さか}をつないでも、\R{急}{きゅう}さは\R{変}{か}わりません。\R{横}{よこ}\qty{20}{m}どうしでくらべると、Aは\qty{2}{m}、Bは\qty{4}{m}\R{上}{あ}がるので、Bのほうが\R{急}{きゅう}です。

\begin{bkfigure}
\begin{tikzpicture}[x=1mm, y=1mm]
  % 坂B（横 25 mm・高さ 5 mm）を4つ、階段のようにつなぐ。1 m を 5 mm でかく
  \foreach \k in {0,...,3} {\bkthrfslope{(\k*25,\k*5)}{25}{5}{bkOrange}
    \node[font=\sffamily\scriptsize, text=bkOrange!80!black, anchor=north, inner sep=1pt] at (\k*25+15,\k*5-0.6) {B};}
  \draw[bkBlue!85!black, line width=1.2pt] (0,0) -- (100,10);
  \draw[bkGray, line width=0.5pt] (0,0) -- (100,0) -- (100,20);
  \node[font=\sffamily\small, text=bkBlue!80!black, anchor=west] at (101.5,10) {A：\qty{2}{m}};
  \node[font=\sffamily\small, text=bkOrange!80!black, anchor=west] at (101.5,20) {B×4：\qty{4}{m}};
  \node[bk tick, anchor=north] at (50,-1.5) {\R{横}{よこ}\qty{20}{m}};
\end{tikzpicture}
\bkfigcaption{\R{坂}{さか}B（\R{横}{よこ}\qty{5}{m}・\R{高}{たか}さ\qty{1}{m}）を4つつなぐと、\R{横}{よこ}\qty{20}{m}で\qty{4}{m}\R{上}{あ}がる\R{坂}{さか}になる。}
\end{bkfigure}

\bkthrfstep{2}{\R{方法}{ほう|ほう}2：\R{横}{よこ}に\qty{1}{m}\R{進}{すす}むと\R{何}{なん}m\R{上}{あ}がるか}
\R{横}{よこ}を\qty{1}{m}にそろえることもできます。「\R{横}{よこ}\qty{1}{m}あたり」の\R{高}{たか}さは、1あたりの\R{量}{りょう}なので、わり\R{算}{ざん}で\R{求}{もと}めます（2-3）。
\bkthrfline{A：$\qty{2}{m}\div\qty{20}{m}=0.1$\qquad B：$\qty{1}{m}\div\qty{5}{m}=0.2$}
Bの0.2はAの0.1の2\R{倍}{ばい}です。\R{方法}{ほう|ほう}1でも、Bの\qty{4}{m}はAの\qty{2}{m}の2\R{倍}{ばい}でした。2つの\R{方法}{ほう|ほう}の\R{答}{こた}えが\R{合}{あ}っています。m÷mなので、この\R{数}{かず}には\R{単位}{たん|い}がつきません。

\bkhatenaref の\R{答}{こた}え：きついのは\R{坂}{さか}Bです。\R{急}{きゅう}さは、Aの2\R{倍}{ばい}あります。

このように、\R{急}{きゅう}さは「\R{上}{あ}がった\R{高}{たか}さ÷\R{横}{よこ}に\R{進}{すす}んだ\R{長}{なが}さ」で\R{表}{あらわ}せます。この\R{数}{かず}を\textbf{\R{傾}{かたむ}き}といいます。\bkidx{かたむき}{\R{傾}{かたむ}き}\R{道路}{どう|ろ}や\R{線路}{せん|ろ}では「こう\R{配}{ばい}」ともよびます。

\bkthrfstep{3}{グラフの\R{傾}{かたむ}き}
グラフの\R{直線}{ちょく|せん}でも\R{同}{おな}じです。\R{直線}{ちょく|せん}の\R{上}{うえ}に2つの\R{点}{てん}をとります。\R{左}{ひだり}の\R{点}{てん}をA、\R{右}{みぎ}の\R{点}{てん}をBとします。AからBへ\R{行}{い}くとき、\R{横}{よこ}に\R{進}{すす}んだ\R{長}{なが}さは $x$ の\R{変化量}{へん|か|りょう} $\Delta x$、\R{上}{あ}がった\R{高}{たか}さは $y$ の\R{変化量}{へん|か|りょう} $\Delta y$ です（どちらも「あと－まえ」。2-3）。

\noindent\begin{minipage}[c]{0.52\linewidth}\centering
\begin{tikzpicture}
\begin{axis}[bkplot, width=68mm, height=56mm, xmin=0, xmax=6.6, ymin=0, ymax=5.4, xtick={1,...,6}, ytick={1,...,5}, xlabel={$x$}, ylabel={$y$}]
  \addplot[bkcurve, domain=0:6.4] {0.5*x+1.5};
  \draw[bkRed, line width=0.9pt] (axis cs:1,2) -- (axis cs:5,2);
  \draw[bkGreen!70!black, line width=0.9pt] (axis cs:5,2) -- (axis cs:5,4);
  \addplot[only marks, mark=*, mark size=2pt, bkInk] coordinates {(1,2) (5,4)};
  \node[anchor=south east, font=\sffamily\small, inner sep=1.5pt] at (axis cs:1,2.05) {A};
  \node[anchor=south east, font=\sffamily\small, inner sep=1.5pt] at (axis cs:5,4.05) {B};
  \node[anchor=north, font=\sffamily\footnotesize, text=bkRed!80!black, fill=white, inner sep=1pt] at (axis cs:3,1.85) {$\Delta x=5-1=4$};
  \node[anchor=west, font=\sffamily\footnotesize, text=bkGreen!50!black, fill=white, inner sep=1pt] at (axis cs:5.12,3) {$\Delta y=4-2=2$};
\end{axis}
\end{tikzpicture}
\end{minipage}\hfill
\begin{minipage}[c]{0.45\linewidth}
A は $(1,\ 2)$、B は $(5,\ 4)$ です。\par\vspace{2pt}
\R{横}{よこ}の\R{変化量}{へん|か|りょう}（\R{赤}{あか}）は、$x$ のあと－まえで $\Delta x=5-1=4$。\par\vspace{2pt}
たての\R{変化量}{へん|か|りょう}（\R{緑}{みどり}）は、$y$ のあと－まえで $\Delta y=4-2=2$。\par\vspace{2pt}
\R{傾}{かたむ}きは $2\div4=\dfrac{2}{4}=0.5$ です。
\end{minipage}

\par\vspace{4pt}
\R{点}{てん}A $(x_1,\ y_1)$ から\R{点}{てん}B $(x_2,\ y_2)$ までなら、$\Delta x=x_2-x_1$、$\Delta y=y_2-y_1$ です。だから、
\bkthrfline{\fbox{\R{傾}{かたむ}き＝たての\R{変化量}{へん|か|りょう}÷\R{横}{よこ}の\R{変化量}{へん|か|りょう}＝$\dfrac{\Delta y}{\Delta x}=\dfrac{y_2-y_1}{x_2-x_1}$}}
と\R{書}{か}けます。$x_1$、$x_2$ の\R{右下}{みぎ|した}の\R{小}{ちい}さな1、2は「1つめの\R{点}{てん}」「2つめの\R{点}{てん}」の\R{印}{しるし}です（3-2）。たてを\R{上}{うえ}に、\R{横}{よこ}を\R{下}{した}に\R{書}{か}きます。\R{坂}{さか}の「\R{高}{たか}さ÷\R{横}{よこ}」と\R{同}{おな}じ\R{順番}{じゅん|ばん}です。

\begin{seikaku}
\textbf{\R{逆}{ぎゃく}の\R{順番}{じゅん|ばん}でひいても\R{同}{おな}じ}\par
BからAへ\R{向}{む}かって\R{計算}{けい|さん}すると、$\Delta y=2-4=-2$、$\Delta x=1-5=-4$ で、$\dfrac{-2}{-4}=0.5$。マイナスどうしのわり\R{算}{ざん}はプラス（2-2）なので、\R{答}{こた}えは\R{同}{おな}じです。ただし、たてと\R{横}{よこ}で\R{同}{おな}じ\R{順番}{じゅん|ばん}にひくこと。たてはBからA、\R{横}{よこ}はAからBのようにまぜると、\R{符号}{ふ|ごう}がまちがいます。
\end{seikaku}

\bkthrfstep{4}{\R{直線}{ちょく|せん}なら、どの2\R{点}{てん}でも\R{傾}{かたむ}きは\R{同}{おな}じ}
\R{同}{おな}じ\R{直線}{ちょく|せん}の\R{上}{うえ}で、\R{別}{べつ}の2\R{点}{てん}をとったら、\R{傾}{かたむ}きは\R{変}{か}わるでしょうか。\R{上}{うえ}の\R{直線}{ちょく|せん}で、Aと $(3,\ 3)$ の\R{点}{てん}Cをとってみます。

\noindent\begin{minipage}[c]{0.52\linewidth}\centering
\begin{tikzpicture}
\begin{axis}[bkplot, width=68mm, height=56mm, xmin=0, xmax=6.6, ymin=0, ymax=5.4, xtick={1,...,6}, ytick={1,...,5}, xlabel={$x$}, ylabel={$y$}]
  \fill[bkOrange!20] (axis cs:1,2) -- (axis cs:5,2) -- (axis cs:5,4) -- cycle;
  \fill[bkBlue!25] (axis cs:1,2) -- (axis cs:3,2) -- (axis cs:3,3) -- cycle;
  \draw[bkBlue!80!black, line width=0.7pt] (axis cs:1,2) -- (axis cs:3,2) -- (axis cs:3,3);
  \draw[bkOrange!80!black, line width=0.7pt] (axis cs:3,2) -- (axis cs:5,2) -- (axis cs:5,4);
  \addplot[bkcurve, domain=0:6.4] {0.5*x+1.5};
  \addplot[only marks, mark=*, mark size=2pt, bkInk] coordinates {(1,2) (3,3) (5,4)};
  \node[anchor=south east, font=\sffamily\small, inner sep=1.5pt] at (axis cs:1,2.05) {A};
  \node[anchor=south east, font=\sffamily\small, inner sep=1.5pt] at (axis cs:3,3.05) {C};
  \node[anchor=south east, font=\sffamily\small, inner sep=1.5pt] at (axis cs:5,4.05) {B};
  \node[anchor=north, font=\sffamily\footnotesize, inner sep=1pt] at (axis cs:2,1.9) {2};
  \node[anchor=west, font=\sffamily\footnotesize, inner sep=1pt] at (axis cs:3.05,2.5) {1};
  \node[anchor=north, font=\sffamily\footnotesize, inner sep=1pt] at (axis cs:4,1.9) {4（AからB）};
  \node[anchor=west, font=\sffamily\footnotesize, inner sep=1pt] at (axis cs:5.08,3) {2};
\end{axis}
\end{tikzpicture}
\end{minipage}\hfill
\begin{minipage}[c]{0.45\linewidth}
AからC：$\Delta x=2$、$\Delta y=1$。\R{傾}{かたむ}きは $\dfrac{1}{2}=0.5$。\par\vspace{2pt}
AからB：$\Delta x=4$、$\Delta y=2$。\R{傾}{かたむ}きは $\dfrac{2}{4}=0.5$。\par\vspace{2pt}
\R{大}{おお}きい\R{三角形}{さん|かく|けい}（だいだい）は、\R{小}{ちい}さい\R{三角形}{さん|かく|けい}（\R{青}{あお}）を、\R{横}{よこ}もたても2\R{倍}{ばい}にのばした\R{形}{かたち}です。
\end{minipage}

\par\vspace{4pt}
\R{横}{よこ}が2\R{倍}{ばい}になると、たても2\R{倍}{ばい}になります。\R{分数}{ぶん|すう}の\R{分母}{ぶん|ぼ}と\R{分子}{ぶん|し}に\R{同}{おな}じ2をかけても、\R{分数}{ぶん|すう}の\R{大}{おお}きさは\R{変}{か}わりません（$\dfrac{1}{2}=\dfrac{2}{4}$。2-4）。まっすぐな\R{線}{せん}は、\R{同}{おな}じ\R{形}{かたち}の\R{階段}{かい|だん}がくり\R{返}{かえ}してできているので（2-15）、どこで\R{何段}{なん|だん}ぶんとっても、たて÷\R{横}{よこ}は\R{同}{おな}じです。だから、\R{直線}{ちょく|せん}の\R{傾}{かたむ}きは、1つに\R{決}{き}まります。

\noindent\begin{minipage}[c]{0.40\linewidth}
1つの\R{直線}{ちょく|せん}に、1つの\R{傾}{かたむ}き。この\R{傾}{かたむ}きを、\R{記号}{き|ごう} $\dfrac{\Delta y}{\Delta x}$ で\R{表}{あらわ}すことがよくあります。
\end{minipage}\hfill
\begin{minipage}[c]{67mm}\bkkigocard{$\dfrac{\Delta y}{\Delta x}$}{デルタエックス\R{分}{ぶん}のデルタワイ}{\R{傾}{かたむ}き。たての\R{変化量}{へん|か|りょう}÷\R{横}{よこ}の\R{変化量}{へん|か|りょう}。}{slope}{3-5}\end{minipage}

\bksubsection{2}{\R{右上}{みぎ|あ}がりは\R{正}{せい}、\R{右下}{みぎ|さ}がりは\R{負}{ふ}、\R{水平}{すい|へい}は0}

\leavevmode\label{bk:3-5:fugou}\R{坂}{さか}には、のぼり\R{坂}{ざか}もくだり\R{坂}{ざか}もあります。グラフも\R{同}{おな}じです。\R{左}{ひだり}から\R{右}{みぎ}へ\R{見}{み}ていくとき、\R{上}{あ}がっていく\R{線}{せん}、\R{下}{さ}がっていく\R{線}{せん}、\R{平}{たい}らな\R{線}{せん}の3つの\R{傾}{かたむ}きを\R{計算}{けい|さん}してみます。どれも、\R{左}{ひだり}の\R{点}{てん}から\R{右}{みぎ}の\R{点}{てん}へ\R{向}{む}かって、あと－まえで\R{計算}{けい|さん}します。

\begin{bkfigure}
\begin{tikzpicture}
\begin{axis}[bkplot, name=g1, width=44mm, height=44mm, xmin=0, xmax=4.6, ymin=0, ymax=4.6, xtick={1,2,3,4}, ytick={1,2,3,4}, xlabel={$x$}, ylabel={$y$}, title={\sffamily\small \R{右上}{みぎ|あ}がり}, title style={yshift=3mm}]
  \addplot[bkcurve, domain=0.3:2.2] {2*x-0.5};
  \draw[bkGray, line width=0.6pt, densely dashed] (axis cs:1,1.5) -- (axis cs:2,1.5) -- (axis cs:2,3.5);
  \addplot[only marks, mark=*, mark size=1.7pt, bkInk] coordinates {(1,1.5) (2,3.5)};
  \node[anchor=west, font=\sffamily\scriptsize, inner sep=1pt] at (axis cs:2.1,2.5) {$+2$};
  \node[anchor=north, font=\sffamily\scriptsize, inner sep=1pt] at (axis cs:1.5,1.4) {$+1$};
\end{axis}
\begin{axis}[bkplot, name=g2, at={($(g1.east)+(9mm,0)$)}, anchor=west, width=44mm, height=44mm, xmin=0, xmax=4.6, ymin=0, ymax=4.6, xtick={1,2,3,4}, ytick={1,2,3,4}, xlabel={$x$}, ylabel={$y$}, title={\sffamily\small \R{右下}{みぎ|さ}がり}, title style={yshift=3mm}]
  \addplot[bkcurve2, domain=0.2:4] {-x+4.2};
  \draw[bkGray, line width=0.6pt, densely dashed] (axis cs:1,3.2) -- (axis cs:3,3.2) -- (axis cs:3,1.2);
  \addplot[only marks, mark=*, mark size=1.7pt, bkInk] coordinates {(1,3.2) (3,1.2)};
  \node[anchor=west, font=\sffamily\scriptsize, inner sep=1pt] at (axis cs:3.1,2.2) {$-2$};
  \node[anchor=south, font=\sffamily\scriptsize, inner sep=1pt] at (axis cs:2,3.3) {$+2$};
\end{axis}
\begin{axis}[bkplot, name=g3, at={($(g2.east)+(9mm,0)$)}, anchor=west, width=44mm, height=44mm, xmin=0, xmax=4.6, ymin=0, ymax=4.6, xtick={1,2,3,4}, ytick={1,2,3,4}, xlabel={$x$}, ylabel={$y$}, title={\sffamily\small \R{水平}{すい|へい}}, title style={yshift=3mm}]
  \addplot[bkcurve3, domain=0.2:4.2] {2};
  \addplot[only marks, mark=*, mark size=1.7pt, bkInk] coordinates {(1,2) (4,2)};
  \node[anchor=south, font=\sffamily\scriptsize, inner sep=1pt] at (axis cs:2.5,2.2) {$+3$};
\end{axis}
\end{tikzpicture}
\bkfigcaption{\R{左}{ひだり}の\R{点}{てん}から\R{右}{みぎ}の\R{点}{てん}へ。\R{横}{よこ}の\R{変化量}{へん|か|りょう}は、どれも\R{正}{せい}（\R{右}{みぎ}へ\R{進}{すす}む）。}
\end{bkfigure}

\begin{bkcenter}
\begin{bktable}{colspec={Q[l,m,wd=22mm] Q[c,m,wd=24mm] Q[c,m,wd=24mm] Q[c,m,wd=48mm]}, rowsep=1.5pt, row{1}={font=\sffamily\small}}
 & $\Delta x$ & $\Delta y$ & \R{傾}{かたむ}き $\dfrac{\Delta y}{\Delta x}$ \\
\R{右上}{みぎ|あ}がり & $+1$ & $+2$ & $\dfrac{2}{1}=2$（\R{正}{せい}） \\
\R{右下}{みぎ|さ}がり & $+2$ & $-2$ & $\dfrac{-2}{2}=-1$（\R{負}{ふ}） \\
\R{水平}{すい|へい} & $+3$ & $0$ & $\dfrac{0}{3}=0$ \\
\end{bktable}
\bktabcaption{\R{右下}{みぎ|さ}がりでは、あと－まえの $\Delta y$ がマイナスになる。}
\end{bkcenter}

\R{右下}{みぎ|さ}がりの\R{線}{せん}では、\R{右}{みぎ}へ\R{進}{すす}むと\R{下}{さ}がるので、$y$ の「あと」が「まえ」より\R{小}{ちい}さくなります。だから $\Delta y$ はマイナスです。マイナスをプラスでわると、\R{答}{こた}えはマイナスです（2-2）。\R{水平}{すい|へい}な\R{線}{せん}は、\R{右}{みぎ}へ\R{進}{すす}んでも\R{高}{たか}さが\R{変}{か}わらないので、$\Delta y=0$ で、\R{傾}{かたむ}きは0です。

\begin{bkcenter}
\begin{bktable}{colspec={Q[c,m,wd=30mm] Q[l,m,wd=84mm]}, rowsep=1.5pt, row{1}={font=\sffamily\bfseries\small}}
\R{傾}{かたむ}き & \R{線}{せん}のようす \\
\R{正}{せい}（＋） & \R{右上}{みぎ|あ}がり。\R{数}{かず}が\R{大}{おお}きいほど\R{急}{きゅう}。 \\
0 & \R{水平}{すい|へい}（\R{平}{たい}ら）。 \\
\R{負}{ふ}（－） & \R{右下}{みぎ|さ}がり。\R{絶対値}{ぜっ|たい|ち}（3-2）が\R{大}{おお}きいほど\R{急}{きゅう}。 \\
\end{bktable}
\bktabcaption{\R{傾}{かたむ}きの\R{符号}{ふ|ごう}は\R{向}{む}き（のぼりか、くだりか）、\R{絶対値}{ぜっ|たい|ち}は\R{急}{きゅう}さを\R{表}{あらわ}す。}
\end{bkcenter}

たとえば、\R{傾}{かたむ}き $-3$ の\R{線}{せん}と\R{傾}{かたむ}き $-1$ の\R{線}{せん}では、\R{数}{かず}としては $-3$ のほうが\R{小}{ちい}さい（2-2）のですが、\R{坂}{さか}としては $-3$ のほうが\R{急}{きゅう}なくだり\R{坂}{ざか}です。\R{右}{みぎ}へ1\R{進}{すす}むと、3も\R{下}{さ}がるからです。\R{急}{きゅう}さをくらべるときは、\R{符号}{ふ|ごう}をとった\R{絶対値}{ぜっ|たい|ち} $|-3|=3$ と $|-1|=1$ をくらべます。

\begin{seikaku}
\textbf{たての\R{直線}{ちょく|せん}には、\R{傾}{かたむ}きがない}\par
$x$ \R{軸}{じく}に\R{垂直}{すい|ちょく}な、まっすぐたての\R{直線}{ちょく|せん}では、2\R{点}{てん}の $x$ が\R{同}{おな}じなので、$\Delta x=0$ です。0でわることはできない（1-14）ので、たての\R{直線}{ちょく|せん}には\R{傾}{かたむ}きの\R{数}{かず}が\R{決}{き}まりません。\R{横}{よこ}に\R{進}{すす}まずに\R{上}{あ}がる、かべのような\R{坂}{さか}だからです。
\end{seikaku}

\bksubsection{3}{まっすぐなグラフの\R{式}{しき}：1\R{次関数}{じ|かん|すう} $y=ax+b$}

\noindent\begin{minipage}[c]{0.66\linewidth}\label{bk:3-5:ichiji}
\R{水}{みず}そうに、はじめから\qty{5}{L}の\R{水}{みず}が\R{入}{はい}っています。じゃ\R{口}{ぐち}から、1\R{分}{ぷん}に\qty{2}{L}ずつ\R{水}{みず}を\R{入}{い}れます。\R{入}{い}れはじめてから $x$ \R{分後}{ふん|ご}の\R{水}{みず}の\R{量}{りょう}を $y$〔L〕とします。3\R{分後}{ぷん|ご}は\R{何}{なん}Lかな？　10\R{分後}{ぷん|ご}は？
\end{minipage}\hfill
\begin{minipage}[c]{0.3\linewidth}\centering
\begin{tikzpicture}[x=1mm, y=1mm]
  \bkthrftank{(4,0)}{9}
  \node[font=\sffamily\scriptsize, anchor=west] at (27,7) {\qty{5}{L}};
  \node[font=\sffamily\scriptsize, anchor=west, text=bkBlue!80!black] at (12,30) {1\R{分}{ぷん}に\qty{2}{L}};
\end{tikzpicture}
\end{minipage}

\par\vspace{4pt}
1\R{分}{ぷん}ごとに\qty{2}{L}ずつふえるので、\R{表}{ひょう}にすると\R{次}{つぎ}のようになります。

\begin{bkcenter}
\begin{bktable}{colspec={Q[l,m,wd=32mm] *{6}{Q[c,m,wd=11mm]}}, rowsep=1.5pt}
\R{時間}{じ|かん} $x$（\R{分}{ふん}） & 0 & 1 & 2 & 3 & 4 & 5 \\
\R{水}{みず}の\R{量}{りょう} $y$（L） & 5 & 7 & 9 & 11 & 13 & 15 \\
\end{bktable}
\bktabcaption{$x$ が1ふえるごとに、$y$ は2ずつふえる。}
\end{bkcenter}

\textbf{\R{式}{しき}にする}\quad $x$ \R{分}{ふん}で\R{入}{はい}る\R{水}{みず}は、\qty{2}{L}の $x$ こ\R{分}{ぶん}で $2\times x$（L）です（1-12：1つ\R{分}{ぶん}×いくつ\R{分}{ぶん}）。はじめの\qty{5}{L}にこれがたされるので、$y=5+2x$。たし\R{算}{ざん}は\R{順番}{じゅん|ばん}を\R{入}{い}れかえてもよいので、$x$ の\R{項}{こう}を\R{先}{さき}に\R{書}{か}いて、
\bkthrfline{$y=2x+5$}
です。$x=3$ なら $y=2\times3+5=11$、$x=10$ なら $y=2\times10+5=25$。3\R{分後}{ぷん|ご}は\qty{11}{L}、10\R{分後}{ぷん|ご}は\qty{25}{L}です。

\noindent\begin{minipage}[c]{0.5\linewidth}\centering
\begin{tikzpicture}
\begin{axis}[bkplot, width=62mm, height=58mm, xmin=0, xmax=5.8, ymin=0, ymax=17.5, xtick={1,...,5}, ytick={5,10,15}, xlabel={$x$（\R{分}{ふん}）}, ylabel={$y$（L）}]
  \addplot[bkcurve, domain=0:5.5] {2*x+5};
  \addplot[bkcurve2, domain=0:5.5] {2*x};
  \addplot[only marks, mark=*, mark size=1.7pt, bkInk] coordinates {(0,5) (1,7) (2,9) (3,11) (4,13) (5,15)};
  \foreach \xx in {1,3} {\edef\tmp{\noexpand\draw[-{Stealth[length=3pt]}, bkGreen!70!black, line width=0.7pt] (axis cs:\xx,{2*\xx+0.3}) -- (axis cs:\xx,{2*\xx+4.7});}\tmp}
  \node[anchor=west, font=\sffamily\scriptsize, text=bkGreen!50!black, inner sep=1pt] at (axis cs:3.08,8.5) {$+5$};
  \node[anchor=east, font=\sffamily\footnotesize, text=bkBlue!80!black, inner sep=1pt] at (axis cs:4.3,15.5) {$y=2x+5$};
  \node[anchor=west, font=\sffamily\footnotesize, text=bkOrange!80!black, inner sep=1pt] at (axis cs:4.3,7) {$y=2x$};
\end{axis}
\end{tikzpicture}
\end{minipage}\hfill
\begin{minipage}[c]{0.47\linewidth}
\R{点}{てん}をかきこむと、まっすぐにならびます。$x$ が1ふえるごとに $y$ がいつも2ずつふえるので、\R{同}{おな}じ\R{形}{かたち}の\R{階段}{かい|だん}がくり\R{返}{かえ}すからです（2-15）。\par\vspace{2pt}
オレンジは、はじめが0Lのとき（$y=2x$、\R{比例}{ひ|れい}）です。\R{青}{あお}の\R{線}{せん}は、オレンジの\R{線}{せん}を、そのまま\R{上}{うえ}へ5ずらしたものです。どの $x$ でも、$y$ が5だけ\R{大}{おお}きいからです（\R{緑}{みどり}の\R{矢印}{や|じるし}）。
\end{minipage}

\par\vspace{4pt}
$2x+5$ のように、「$x$ に\R{数}{かず}をかけた\R{項}{こう}」と「\R{数}{かず}だけの\R{項}{こう}」でできた\R{式}{しき}を、$x$ の\textbf{1\R{次式}{じ|しき}}といいます（$x^2$ のような2\R{乗}{じょう}の\R{項}{こう}がないから「1\R{次}{じ}」）。$y$ が $x$ の1\R{次式}{じ|しき}で\R{書}{か}けるとき、「$y$ は $x$ の\textbf{1\R{次関数}{じ|かん|すう}}である」といいます。\bkidx{いちじかんすう}{1\R{次関数}{じ|かん|すう}}\R{一般}{いっ|ぱん}に、1\R{次関数}{じ|かん|すう}は
\bkthrfline{\fbox{$y=ax+b$}\qquad（$a$、$b$ は\R{定数}{てい|すう}。$a$ は0でない）}
と\R{書}{か}けます。\R{定数}{てい|すう}は「$x$ が\R{変}{か}わっても\R{変}{か}わらない\R{数}{かず}」でした（3-1）。\R{水}{みず}そうなら $a=2$、$b=5$ です。\R{比例}{ひ|れい} $y=ax$ は、$b=0$ の1\R{次関数}{じ|かん|すう}です。

\bkthrfstep{1}{$a$ は\R{傾}{かたむ}き：\R{変化}{へん|か}の\R{割合}{わり|あい}}
\R{表}{ひょう}で、$x$ が1から4まで\R{変}{か}わるときを\R{見}{み}ます。$y$ は7から13まで\R{変}{か}わります。
\[ \frac{\Delta y}{\Delta x}=\frac{13-7}{4-1}=\frac{6}{3}=2 \]
$x$ が0から5までなら、$\dfrac{15-5}{5-0}=\dfrac{10}{5}=2$。どこを\R{選}{えら}んでも2です。この2は、「1\R{分}{ぷん}あたり\qty{2}{L}」の2、つまり $y=2x+5$ の $a$ です。

$x$ の\R{変化量}{へん|か|りょう}に\R{対}{たい}する $y$ の\R{変化量}{へん|か|りょう}の\R{割合}{わり|あい} $\dfrac{\Delta y}{\Delta x}$ を、\textbf{\R{変化}{へん|か}の\R{割合}{わり|あい}}といいます。\bkidx{へんかのわりあい}{\R{変化}{へん|か}の\R{割合}{わり|あい}}\R{中学校}{ちゅう|がっ|こう}の\R{教科書}{きょう|か|しょ}では、\R{変化量}{へん|か|りょう}を「\R{増加量}{ぞう|か|りょう}」とよぶこともあります。1\R{次関数}{じ|かん|すう}では、\R{変化}{へん|か}の\R{割合}{わり|あい}はグラフの\R{傾}{かたむ}きと\R{同}{おな}じで、いつも $a$ です。

\textbf{なぜいつも $a$ になるの？}\quad \R{文字}{も|じ}のままで\R{確}{たし}かめます。$x$ が $x_1$ から $x_2$ まで\R{変}{か}わるとき、$y_1=ax_1+b$、$y_2=ax_2+b$ です。
\begin{align*}
\Delta y=y_2-y_1&=(ax_2+b)-(ax_1+b)\\
&=ax_2+b-ax_1-b && \text{\small（かっこをはずす。2-12）}\\
&=ax_2-ax_1=a(x_2-x_1) && \text{\small（$b-b=0$。$a$ でくくる）}
\end{align*}
$x_2-x_1$ は $\Delta x$ なので、$\Delta y=a\times\Delta x$。\R{両辺}{りょう|へん}を $\Delta x$ でわると $\dfrac{\Delta y}{\Delta x}=a$ です。はじめの\R{量}{りょう} $b$ は、あと－まえのひき\R{算}{ざん}で\R{消}{き}えてしまいます。\R{変化}{へん|か}の\R{速}{はや}さは、はじめにいくら\R{入}{はい}っていたかとは\R{関係}{かん|けい}がないのです。

\bkthrfstep{2}{$b$ は\R{切片}{せっ|ぺん}：グラフが $y$ \R{軸}{じく}と\R{交}{まじ}わる\R{高}{たか}さ}
$x=0$ を\R{入}{い}れると、$y=a\times0+b=b$ です。グラフでは、$x=0$ の\R{所}{ところ}、つまり $y$ \R{軸}{じく}の\R{上}{うえ}の\R{点}{てん} $(0,\ b)$ を\R{通}{とお}ります。この $b$ を、グラフの\textbf{\R{切片}{せっ|ぺん}}といいます。\bkidx{せっぺん}{\R{切片}{せっ|ぺん}}$y$ \R{軸}{じく}から\R{切}{き}りとる\R{所}{ところ}の\R{高}{たか}さ、という\R{意味}{い|み}です。\R{水}{みず}そうなら、\R{入}{い}れはじめ（0\R{分}{ふん}）の\qty{5}{L}が\R{切片}{せっ|ぺん}です。

\bkthrfstep{3}{$a$ が\R{負}{ふ}のとき}
スマートフォンの\R{電池}{でん|ち}が100\%（パーセント）あります。\R{動画}{どう|が}を\R{見}{み}つづけると、1\R{時間}{じ|かん}に20\%ずつへります。$x$ \R{時間後}{じ|かん|ご}の\R{電池}{でん|ち}の\R{残}{のこ}りを $y$\,\% とします。

\begin{bkcenter}
\begin{bktable}{colspec={Q[l,m,wd=30mm] *{6}{Q[c,m,wd=11mm]}}, rowsep=1pt}
\R{時間}{じ|かん} $x$（h） & 0 & 1 & 2 & 3 & 4 & 5 \\
\R{残}{のこ}り $y$（\%） & 100 & 80 & 60 & 40 & 20 & 0 \\
\end{bktable}
\bktabcaption{h は\R{時間}{じ|かん}の\R{記号}{き|ごう}（hour の h。2-10）。}
\end{bkcenter}
\noindent\begin{minipage}[c]{0.54\linewidth}
1\R{時間}{じ|かん}に20ずつへるので、$x$ \R{時間}{じ|かん}では $20x$ へります。はじめの100から\R{引}{ひ}いて $y=100-20x$。\R{並}{なら}べかえると
\bkthrfline{$y=-20x+100$}
です。$a=-20$、$b=100$。\R{傾}{かたむ}きが\R{負}{ふ}なので、グラフは\R{右下}{みぎ|さ}がりです。$-20$ は「1\R{時間}{じ|かん}あたり20\%へる」という、へる\R{速}{はや}さを\R{表}{あらわ}しています。
\end{minipage}\hfill
\begin{minipage}[c]{0.43\linewidth}\centering
\begin{tikzpicture}
\begin{axis}[bkplot, width=48mm, height=48mm, xmin=0, xmax=5.8, ymin=0, ymax=115, xtick={1,...,5}, ytick={20,40,...,100}, xlabel={$x$（h）}, ylabel={$y$（\%）}]
  \addplot[bkcurve2, domain=0:5] {-20*x+100};
  \addplot[only marks, mark=*, mark size=1.6pt, bkInk] coordinates {(0,100) (1,80) (2,60) (3,40) (4,20) (5,0)};
  \node[anchor=south west, font=\sffamily\footnotesize, text=bkOrange!80!black, inner sep=1pt] at (axis cs:2.6,58) {$y=-20x+100$};
\end{axis}
\end{tikzpicture}
\end{minipage}

\bkthrfstep{4}{\R{式}{しき}からグラフをかく\R{手順}{て|じゅん}}
1\R{次関数}{じ|かん|すう}のグラフは\R{直線}{ちょく|せん}なので、2つの\R{点}{てん}がわかればかけます。$y=\dfrac{1}{2}x+1$ でやってみましょう。

\noindent\begin{minipage}[c]{0.55\linewidth}
\begin{enumerate}[label={\bkcircled[bkOrange]{\arabic*}}, leftmargin=2.4\zw, itemsep=2pt, topsep=0pt]
\item \R{切片}{せっ|ぺん}の\R{点}{てん}をとる。$b=1$ なので $(0,\ 1)$。
\item そこから\R{傾}{かたむ}きの\R{分}{ぶん}だけ\R{進}{すす}む。\R{傾}{かたむ}き $\dfrac{1}{2}$ は「\R{右}{みぎ}へ1で\R{上}{うえ}へ $\dfrac{1}{2}$」。\R{目}{め}もりにのらないので、2\R{倍}{ばい}して「\R{右}{みぎ}へ2で\R{上}{うえ}へ1」と\R{進}{すす}む（\R{傾}{かたむ}きは $\dfrac{1}{2}=\dfrac{2}{4}=\cdots$ のどれでも\R{同}{おな}じ）。$(2,\ 2)$ に\R{着}{つ}く。
\item 2つの\R{点}{てん}を\R{通}{とお}る\R{直線}{ちょく|せん}を\R{引}{ひ}く。\R{傾}{かたむ}きが\R{負}{ふ}のときは、\R{右}{みぎ}へ\R{進}{すす}んで\R{下}{した}へ\R{下}{さ}がる。
\end{enumerate}
\end{minipage}\hfill
\begin{minipage}[c]{0.42\linewidth}\centering
\begin{tikzpicture}
\begin{axis}[bkplot, width=52mm, height=46mm, xmin=-1.5, xmax=4.6, ymin=-0.8, ymax=3.8, xtick={-1,1,2,3,4}, ytick={1,2,3}, xlabel={$x$}, ylabel={$y$}]
  \addplot[bkcurve, domain=-1.4:4.4] {0.5*x+1};
  \draw[bkRed, line width=0.8pt, -{Stealth[length=3pt]}] (axis cs:0.1,1) -- (axis cs:1.95,1);
  \draw[bkGreen!70!black, line width=0.8pt, -{Stealth[length=3pt]}] (axis cs:2,1.05) -- (axis cs:2,1.95);
  \addplot[only marks, mark=*, mark size=1.8pt, bkInk] coordinates {(0,1) (2,2)};
  \node[anchor=north, font=\sffamily\scriptsize, text=bkRed!80!black, inner sep=1pt] at (axis cs:1.15,0.85) {\R{右}{みぎ}へ2};
  \node[anchor=west, font=\sffamily\scriptsize, text=bkGreen!50!black, inner sep=1pt] at (axis cs:2.1,1.5) {\R{上}{うえ}へ1};
\end{axis}
\end{tikzpicture}
\end{minipage}

\bksubsection{4}{$x$-$t$ グラフの\R{傾}{かたむ}きは、\R{速}{はや}さ}

\leavevmode\label{bk:3-5:hayasa}ミナさんは\R{自転車}{じ|てん|しゃ}で、\R{家}{いえ}から\qty{1000}{m}はなれた\R{公園}{こう|えん}へ\R{行}{い}き、しばらく\R{休}{やす}んでから\R{家}{いえ}へ\R{帰}{かえ}りました。\R{家}{いえ}を0として、\R{公園}{こう|えん}の\R{向}{む}きに\R{何}{なん}mの\R{所}{ところ}にいるかを\R{位置}{い|ち} $x$ とし、\R{出発}{しゅっ|ぱつ}してからの\R{時間}{じ|かん}を $t$ として、$x$-$t$ グラフにしました（3-4）。

\begin{bkfigure}
\begin{tikzpicture}
\begin{axis}[bkplot, width=118mm, height=58mm, xmin=0, xmax=860, ymin=0, ymax=1250, xtick={100,200,...,800}, ytick={200,400,...,1000}, xlabel={$t$（s）}, ylabel={$x$（m）}]
  \addplot[bkcurve, line width=1.8pt] coordinates {(0,0) (200,1000)};
  \addplot[bkcurve3, line width=1.8pt] coordinates {(200,1000) (500,1000)};
  \addplot[bkcurve2, line width=1.8pt] coordinates {(500,1000) (750,0)};
  \addplot[only marks, mark=*, mark size=1.6pt, bkInk] coordinates {(0,0) (200,1000) (500,1000) (750,0)};
  \node[anchor=west, font=\sffamily\footnotesize, text=bkBlue!80!black, fill=white, inner sep=1pt] at (axis cs:118,450) {\R{行}{い}き};
  \node[anchor=south, font=\sffamily\footnotesize, text=bkGreen!50!black, inner sep=1pt] at (axis cs:350,1040) {\R{公園}{こう|えん}で\R{休}{やす}む};
  \node[anchor=west, font=\sffamily\footnotesize, text=bkOrange!80!black, fill=white, inner sep=1pt] at (axis cs:655,450) {\R{帰}{かえ}り};
\end{axis}
\end{tikzpicture}
\bkfigcaption{ミナさんの $x$-$t$ グラフ。\R{行}{い}きは0〜200\R{秒}{びょう}、\R{休}{やす}みは200〜500\R{秒}{びょう}、\R{帰}{かえ}りは500〜750\R{秒}{びょう}。}\label{bk:3-5:fig-mina}
\end{bkfigure}

3つの\R{部分}{ぶ|ぶん}の\R{傾}{かたむ}きを\R{計算}{けい|さん}します。たては\R{位置}{い|ち} $x$、\R{横}{よこ}は\R{時間}{じ|かん} $t$ なので、\R{傾}{かたむ}きは $\dfrac{\Delta x}{\Delta t}$ です。
\begin{align*}
\text{\small\R{行}{い}き}\quad&\frac{\Delta x}{\Delta t}=\frac{\qty{1000}{m}-\qty{0}{m}}{\qty{200}{s}-\qty{0}{s}}=\frac{\qty{1000}{m}}{\qty{200}{s}}=\qty{5}{m/s}\\
\text{\small\R{休}{やす}み}\quad&\frac{\Delta x}{\Delta t}=\frac{\qty{1000}{m}-\qty{1000}{m}}{\qty{500}{s}-\qty{200}{s}}=\frac{\qty{0}{m}}{\qty{300}{s}}=\qty{0}{m/s}\\
\text{\small\R{帰}{かえ}り}\quad&\frac{\Delta x}{\Delta t}=\frac{\qty{0}{m}-\qty{1000}{m}}{\qty{750}{s}-\qty{500}{s}}=\frac{\qty{-1000}{m}}{\qty{250}{s}}=\qty{-4}{m/s}
\end{align*}

\textbf{\R{傾}{かたむ}き＝\R{速}{はや}さ}\quad $\Delta x\div\Delta t$ は、\R{進}{すす}んだ\R{道}{みち}のり÷かかった\R{時間}{じ|かん}、つまり\R{速}{はや}さそのものです（2-3）。\R{行}{い}きは\R{秒速}{びょう|そく}\qty{5}{m}で\R{走}{はし}っていました。$x$-$t$ グラフでは、\R{線}{せん}が\R{急}{きゅう}なほど、\R{同}{おな}じ\R{時間}{じ|かん}にたくさん\R{進}{すす}むので、\R{速}{はや}いのです。

\textbf{\R{傾}{かたむ}き0＝\R{止}{と}まっている}\quad \R{休}{やす}んでいるあいだは、\R{時間}{じ|かん}がたっても\R{位置}{い|ち}が\R{変}{か}わりません。グラフは\R{水平}{すい|へい}で、\R{速}{はや}さは0です。

\textbf{\R{傾}{かたむ}きが\R{負}{ふ}＝\R{反対}{はん|たい}\R{向}{む}きに\R{進}{すす}む}\quad \R{帰}{かえ}りの\R{傾}{かたむ}き $-4$ のマイナスは、$x$ が\R{小}{ちい}さくなる\R{向}{む}き、つまり\R{家}{いえ}のほうへ\R{進}{すす}んでいることを\R{表}{あらわ}します。\R{速}{はや}さの\R{大}{おお}きさは\R{秒速}{びょう|そく}\qty{4}{m}です。\R{帰}{かえ}りのほうが、\R{行}{い}きより\R{少}{すこ}しゆっくりでした。\R{物理}{ぶつ|り}では、このように\R{向}{む}きまでふくめた\R{速}{はや}さを\textbf{\R{速度}{そく|ど}}とよびます（5-5）。

\textbf{\R{傾}{かたむ}きには\R{単位}{たん|い}がある}\quad \R{坂}{さか}の\R{傾}{かたむ}きは m÷m で\R{単位}{たん|い}がありませんでした。$x$-$t$ グラフの\R{傾}{かたむ}きは、m÷s で m/s という\R{単位}{たん|い}がつきます。\R{傾}{かたむ}きの\R{単位}{たん|い}は、いつも「たての\R{単位}{たん|い}÷\R{横}{よこ}の\R{単位}{たん|い}」です。

\begin{bkcenter}
\begin{bktable}{colspec={Q[l,m,wd=34mm] Q[c,m,wd=20mm] Q[c,m,wd=16mm] Q[c,m,wd=18mm] Q[l,m,wd=30mm]}, rowsep=1pt, cells={font=\small}, row{1}={font=\sffamily\bfseries\small}}
グラフ & たて & \R{横}{よこ} & \R{傾}{かたむ}き & \R{意味}{い|み} \\
\R{位置}{い|ち}と\R{時間}{じ|かん} & m & s & m/s & \R{速}{はや}さ \\
\R{水}{みず}の\R{量}{りょう}と\R{時間}{じ|かん} & L & min & L/min & \R{水}{みず}の\R{入}{はい}る\R{速}{はや}さ \\
\R{電池}{でん|ち}の\R{残}{のこ}りと\R{時間}{じ|かん} & \% & h & \%/h & \R{電池}{でん|ち}のへる\R{速}{はや}さ \\
\R{気温}{き|おん}と\R{時刻}{じ|こく} & \unit{\degreeCelsius} & h & \unit{\degreeCelsius/h} & \R{気温}{き|おん}の\R{上}{あ}がる\R{速}{はや}さ \\
\R{力}{ちから}とばねののび & N & m & N/m & ばね\R{定数}{てい|すう}（2-15） \\
\end{bktable}
\bktabcaption{\R{傾}{かたむ}きの\R{単位}{たん|い}は「たての\R{単位}{たん|い}÷\R{横}{よこ}の\R{単位}{たん|い}」。min は\R{分}{ふん}の\R{記号}{き|ごう}。}
\end{bkcenter}

\R{横}{よこ}が\R{時間}{じ|かん}のグラフでは、\R{傾}{かたむ}きは「1\R{秒}{びょう}あたり・1\R{分}{ぷん}あたり・1\R{時間}{じ|かん}あたり、たての\R{量}{りょう}がどれだけ\R{変}{か}わるか」、つまり\textbf{\R{変化}{へん|か}の\R{速}{はや}さ}です。この\R{章}{しょう}の\R{題}{だい}「グラフの\R{傾}{かたむ}きは\R{変化}{へん|か}の\R{速}{はや}さ」は、このことです。

\bksubsection{5}{2\R{点}{てん}を\R{通}{とお}る\R{直線}{ちょく|せん}の\R{式}{しき}}

\noindent\begin{minipage}[c]{0.6\linewidth}\label{bk:3-5:niten}
ばねにおもりをつるして、ばね\R{全体}{ぜん|たい}の\R{長}{なが}さをはかりました。おもり2\R{個}{こ}で\qty{14}{cm}、5\R{個}{こ}で\qty{20}{cm}でした。ばね\R{全体}{ぜん|たい}の\R{長}{なが}さは、おもりの\R{数}{かず}に\R{比例}{ひ|れい}しません（2-15の「まちがえやすいところ」）。でも、のびは\R{比例}{ひ|れい}するので、グラフは\R{直線}{ちょく|せん}になり、1\R{次関数}{じ|かん|すう}で\R{書}{か}けるはずです。おもりの\R{数}{かず}を $x$ \R{個}{こ}、ばね\R{全体}{ぜん|たい}の\R{長}{なが}さを $y$〔cm〕として、$y=ax+b$ の $a$ と $b$ を\R{求}{もと}めます。
\end{minipage}\hfill
\begin{minipage}[c]{0.38\linewidth}\centering
\begin{tikzpicture}
\begin{axis}[bkplot, width=42mm, height=50mm, xmin=0, xmax=6.4, ymin=0, ymax=24, xtick={1,...,6}, ytick={5,10,15,20}, xlabel={$x$（\R{個}{こ}）}, ylabel={$y$（cm）}]
  \addplot[bkcurve, domain=0:6.2] {2*x+10};
  \addplot[only marks, mark=*, mark size=1.8pt, bkInk] coordinates {(2,14) (5,20)};
  \addplot[only marks, mark=o, mark size=1.8pt, bkRed, line width=0.8pt] coordinates {(0,10)};
  \node[anchor=north west, font=\sffamily\scriptsize, inner sep=1pt] at (axis cs:2,13.6) {$(2,\ 14)$};
  \node[anchor=north west, font=\sffamily\scriptsize, inner sep=1pt] at (axis cs:5,19.6) {$(5,\ 20)$};
  \node[anchor=west, font=\sffamily\scriptsize, text=bkRed!80!black, inner sep=1pt] at (axis cs:0.25,8.6) {\R{切片}{せっ|ぺん}？};
\end{axis}
\end{tikzpicture}
\end{minipage}

\bkthrfstep{1}{\R{傾}{かたむ}き $a$ を\R{求}{もと}める}
2つの\R{点}{てん} $(2,\ 14)$ と $(5,\ 20)$ から、
\[ a=\frac{\Delta y}{\Delta x}=\frac{20-14}{5-2}=\frac{6}{3}=2 \]
おもりを1\R{個}{こ}ふやすごとに、\qty{2}{cm}のびるということです。

\bkthrfstep{2}{1つの\R{点}{てん}を\R{代入}{だい|にゅう}して $b$ を\R{求}{もと}める}
$y=2x+b$ のグラフは $(2,\ 14)$ を\R{通}{とお}るので、$x=2$、$y=14$ を\R{代入}{だい|にゅう}すると\R{等式}{とう|しき}が\R{成}{な}り\R{立}{た}ちます（2-13）。
\begin{align*}
14&=2\times2+b && \text{\small（$x=2$、$y=14$ を\R{代入}{だい|にゅう}）}\\
14&=4+b\\
14-4&=b && \text{\small（4を\R{移項}{い|こう}。2-14）}\\
b&=10
\end{align*}

\bkthrfstep{3}{\R{式}{しき}を\R{書}{か}き、もう1つの\R{点}{てん}で\R{確}{たし}かめる}
$y=2x+10$ です。$x=5$ を\R{入}{い}れると $y=2\times5+10=20$ で、もう1つの\R{点}{てん} $(5,\ 20)$ も\R{通}{とお}っています。\R{切片}{せっ|ぺん}10は、おもり0\R{個}{こ}のときのばねの\R{長}{なが}さ（もとの\R{長}{なが}さ）\qty{10}{cm}です。はかっていない「0\R{個}{こ}」の\R{長}{なが}さが、\R{計算}{けい|さん}でわかりました。

\noindent この\R{求}{もと}め\R{方}{かた}のほかに、$a$ と $b$ を2つの\R{式}{しき}からいっしょに\R{求}{もと}める\R{方法}{ほう|ほう}もあります（3-6）。

\bksubsection{6}{1\R{点}{てん}と\R{傾}{かたむ}きから：$y-y_1=m(x-x_1)$}

\leavevmode\label{bk:3-5:tenkatamuki}こんどは、「おもり2\R{個}{こ}で\qty{14}{cm}。1\R{個}{こ}ふやすごとに\qty{2}{cm}のびる」とだけわかっているとします。\R{通}{とお}る1\R{点}{てん} $(2,\ 14)$ と、\R{傾}{かたむ}き2がわかっています。このとき、\R{直線}{ちょく|せん}の\R{式}{しき}をすぐに\R{書}{か}く\R{方法}{ほう|ほう}があります。

\bkthrfstep{1}{\R{見方}{み|かた}1：\R{傾}{かたむ}きの\R{式}{しき}から}
\noindent\begin{minipage}[c]{0.5\linewidth}\centering
\begin{tikzpicture}
\begin{axis}[bkplot, width=62mm, height=52mm, xmin=0, xmax=4.4, ymin=0, ymax=6.6, xtick={1,3}, xticklabels={$x_1$, $x$}, ytick={2,5}, yticklabels={$y_1$, $y$}, xlabel={}, ylabel={}, grid=none]
  \addplot[bkcurve, domain=0.2:4] {2+1.5*(x-1)};
  \draw[bkRed, line width=0.8pt] (axis cs:1,2) -- (axis cs:3,2);
  \draw[bkGreen!70!black, line width=0.8pt] (axis cs:3,2) -- (axis cs:3,5);
  \draw[bkGray, line width=0.5pt, densely dashed] (axis cs:1,0) -- (axis cs:1,2) -- (axis cs:0,2);
  \draw[bkGray, line width=0.5pt, densely dashed] (axis cs:3,0) -- (axis cs:3,2) (axis cs:3,5) -- (axis cs:0,5);
  \addplot[only marks, mark=*, mark size=1.8pt, bkInk] coordinates {(1,2) (3,5)};
  \node[anchor=south east, font=\sffamily\footnotesize, inner sep=1.5pt] at (axis cs:1,2.1) {A};
  \node[anchor=south east, font=\sffamily\footnotesize, inner sep=1.5pt] at (axis cs:3,5.1) {P};
  \node[anchor=north, font=\sffamily\footnotesize, text=bkRed!80!black, inner sep=1pt] at (axis cs:2,1.85) {$x-x_1$};
  \node[anchor=west, font=\sffamily\footnotesize, text=bkGreen!50!black, inner sep=1pt] at (axis cs:3.1,3.5) {$y-y_1$};
  \node[anchor=south east, font=\sffamily\footnotesize, text=bkBlue!80!black, inner sep=1pt] at (axis cs:2.3,4.4) {\R{傾}{かたむ}き $m$};
\end{axis}
\end{tikzpicture}
\end{minipage}\hfill
\begin{minipage}[c]{0.47\linewidth}
\R{傾}{かたむ}きを $m$ とし、\R{直線}{ちょく|せん}が\R{点}{てん}A $(x_1,\ y_1)$ を\R{通}{とお}るとします。\R{直線}{ちょく|せん}の\R{上}{うえ}に、A でない\R{点}{てん} P $(x,\ y)$ をどこでもいいからとります。AからPまでの\R{傾}{かたむ}きも $m$ です（\R{直線}{ちょく|せん}の\R{傾}{かたむ}きは1つ）。
\end{minipage}

\par\vspace{4pt}
AからPまでは、\R{横}{よこ}の\R{変化量}{へん|か|りょう}が $x-x_1$、たての\R{変化量}{へん|か|りょう}が $y-y_1$ なので、
\[ \frac{y-y_1}{x-x_1}=m \]
\R{両辺}{りょう|へん}に $x-x_1$ をかけると（2-14。\R{左辺}{さ|へん}の\R{分母}{ぶん|ぼ}が\R{約分}{やく|ぶん}されて\R{消}{き}える）、
\bkthrfline{\fbox{$y-y_1=m(x-x_1)$}}
これが、\textbf{\R{傾}{かたむ}きが $m$ で、\R{点}{てん} $(x_1,\ y_1)$ を\R{通}{とお}る\R{直線}{ちょく|せん}の\R{式}{しき}}です。\R{傾}{かたむ}きを $m$ と\R{書}{か}くことも\R{多}{おお}いのですが、$m$ は $y=ax+b$ の $a$ と\R{同}{おな}じもの（\R{傾}{かたむ}き）です。\R{文字}{も|じ}がちがうだけです。

\bkthrfstep{2}{\R{見方}{み|かた}2：\R{階段}{かい|だん}をのぼる}
\R{傾}{かたむ}き $m$ は「\R{右}{みぎ}へ1\R{進}{すす}むと\R{上}{うえ}へ $m$」でした。だから、Aから\R{右}{みぎ}へ $x-x_1$ だけ\R{進}{すす}むと、\R{上}{うえ}へ $m\times(x-x_1)$ だけ\R{上}{あ}がります。Aの\R{高}{たか}さは $y_1$ なので、Pの\R{高}{たか}さは
\bkthrfline{$y=y_1+m(x-x_1)$}
です。$y_1$ を\R{左辺}{さ|へん}へ\R{移項}{い|こう}すると $y-y_1=m(x-x_1)$。\R{見方}{み|かた}1と\R{同}{おな}じ\R{式}{しき}になりました。

\bkthrfstep{3}{\R{数}{かず}で\R{確}{たし}かめる}
ばねなら、$m=2$、$(x_1,\ y_1)=(2,\ 14)$ なので、
\begin{align*}
y-14&=2(x-2) && \text{\small（$m$、$x_1$、$y_1$ に\R{数}{かず}を\R{入}{い}れる）}\\
y-14&=2x-4 && \text{\small（\R{右辺}{う|へん}のかっこをはずす。2-12）}\\
y&=2x-4+14 && \text{\small（$-14$ を\R{移項}{い|こう}。2-14）}\\
y&=2x+10
\end{align*}
2\R{点}{てん}から\R{求}{もと}めた\R{式}{しき}（\pageref{bk:3-5:niten}ページ）と\R{同}{おな}じです。\R{展開}{てん|かい}すると、いつでも $y=ax+b$ の\R{形}{かたち}にもどせます。

\textbf{なぜ\R{分数}{ぶん|すう}の\R{形}{かたち}のままにしないの？}\quad $\dfrac{y-y_1}{x-x_1}=m$ の\R{形}{かたち}では、$x=x_1$ のとき\R{分母}{ぶん|ぼ}が0になって、\R{計算}{けい|さん}できません。つまり、\R{点}{てん}Aそのものが\R{式}{しき}から\R{落}{お}ちてしまいます。かけ\R{算}{ざん}の\R{形}{かたち} $y-y_1=m(x-x_1)$ なら、$x=x_1$ を\R{入}{い}れると $y-y_1=m\times0=0$ で $y=y_1$。\R{点}{てん}Aもちゃんと\R{入}{はい}ります。

\textbf{マイナスの\R{座標}{ざ|ひょう}に\R{注意}{ちゅう|い}}\quad \R{点}{てん} $(-1,\ 3)$ を\R{通}{とお}り、\R{傾}{かたむ}き2の\R{直線}{ちょく|せん}なら、$x_1=-1$ なので $x-x_1=x-(-1)=x+1$ です。
\[ y-3=2(x+1)\quad\to\quad y-3=2x+2\quad\to\quad y=2x+5 \]
$x_1$ がマイナスのときは、かっこをつけてから\R{入}{い}れます（2-13の「かっこごと\R{入}{い}れる」）。

\begin{motto}[3]
\textbf{\R{第}{だい}5\R{巻}{かん}では、この\R{式}{しき}で\R{接線}{せっ|せん}を\R{引}{ひ}く}\par
\R{曲}{ま}がったグラフの1\R{点}{てん}に、ぴったりふれる\R{直線}{ちょく|せん}を\textbf{\R{接線}{せっ|せん}}といいます（5-3）。\R{接線}{せっ|せん}は、ふれる\R{点}{てん}は1つわかっていて、\R{傾}{かたむ}きは\R{微分}{び|ぶん}という\R{計算}{けい|さん}で\R{求}{もと}めます。「1\R{点}{てん}と\R{傾}{かたむ}き」がわかるので、$y-y_1=m(x-x_1)$ がそのまま\R{使}{つか}えます。
\end{motto}

\bksubsection{7}{$v$-$t$ グラフでは、\R{面積}{めん|せき}が\R{道}{みち}のり}

\leavevmode\label{bk:3-5:vt}ミナさんの\R{行}{い}き（\R{図}{ず}\ref{bk:3-5:fig-mina}の0〜200\R{秒}{びょう}。\pageref{bk:3-5:fig-mina}ページ）を、2つのグラフでならべてみます。\R{左}{ひだり}は $x$-$t$ グラフ、\R{右}{みぎ}は\R{横}{よこ}に\R{時間}{じ|かん}、たてに\R{速}{はや}さをとった $v$-$t$ グラフ（2-3）です。

\begin{bkfigure}
\begin{tikzpicture}
\begin{axis}[bkplot, name=xt, width=60mm, height=50mm, xmin=0, xmax=250, ymin=0, ymax=1250, xtick={100,200}, ytick={500,1000}, xlabel={$t$（s）}, ylabel={$x$（m）}]
  \addplot[bkcurve, domain=0:200] {5*x};
  \draw[bkRed, line width=0.8pt] (axis cs:0,0) -- (axis cs:200,0);
  \draw[bkGreen!70!black, line width=0.8pt] (axis cs:200,0) -- (axis cs:200,1000);
  \node[anchor=east, font=\sffamily\footnotesize, text=bkBlue!80!black, fill=white, inner sep=1pt] at (axis cs:95,700) {\R{傾}{かたむ}き\qty{5}{m/s}};
  \node[anchor=west, font=\sffamily\footnotesize, text=bkGreen!50!black, inner sep=1pt] at (axis cs:204,550) {\qty{1000}{m}};
\end{axis}
\begin{axis}[bkplot, at={($(xt.east)+(14mm,0)$)}, anchor=west, width=60mm, height=50mm, xmin=0, xmax=250, ymin=0, ymax=7, xtick={100,200}, ytick={1,...,6}, xlabel={$t$（s）}, ylabel={$v$（m/s）}]
  \fill[bkBlue!20] (axis cs:0,0) rectangle (axis cs:200,5);
  \addplot[bkcurve, domain=0:200] {5};
  \node[font=\sffamily\footnotesize, text=bkBlue!80!black, align=center] at (axis cs:100,2.5) {\R{面積}{めん|せき}\\$5\times200=1000$};
\end{axis}
\end{tikzpicture}
\bkfigcaption{\R{左}{ひだり}の\R{傾}{かたむ}き5が、\R{右}{みぎ}の\R{高}{たか}さ5になる。\R{右}{みぎ}の\R{面積}{めん|せき}1000が、\R{左}{ひだり}の\R{上}{あ}がった\R{高}{たか}さ1000になる。}
\end{bkfigure}

$x$-$t$ グラフの\R{傾}{かたむ}き\qty{5}{m/s}が、$v$-$t$ グラフでは\R{高}{たか}さ5になっています。\R{反対}{はん|たい}に、$v$-$t$ グラフの\R{長方形}{ちょう|ほう|けい}の\R{面積}{めん|せき}（たて×\R{横}{よこ}＝\R{速}{はや}さ×\R{時間}{じ|かん}）は $5\times200=1000$ で、\R{進}{すす}んだ\R{道}{みち}のり\qty{1000}{m}です（2-3）。これは、$x$-$t$ グラフで\R{上}{あ}がった\R{高}{たか}さ $\Delta x$ と\R{同}{おな}じです。
\bkthrfline{$x$-$t$ グラフの\R{傾}{かたむ}き\ →\ $v$-$t$ グラフの\R{高}{たか}さ\qquad $v$-$t$ グラフの\R{面積}{めん|せき}\ →\ $x$ の\R{変化量}{へん|か|りょう}}
\R{速}{はや}さがとちゅうで\R{変}{か}わると、$x$-$t$ グラフは\R{曲}{ま}がり、$v$-$t$ グラフは\R{長方形}{ちょう|ほう|けい}でなくなります。そのときの\R{傾}{かたむ}きは5-1・5-3で、\R{面積}{めん|せき}は5-9で\R{考}{かんが}えます。この「\R{傾}{かたむ}き」と「\R{面積}{めん|せき}」の\R{行}{い}き\R{来}{き}が、\R{第}{だい}5\R{巻}{かん}の\R{微分}{び|ぶん}と\R{積分}{せき|ぶん}です。

\bksubsection{8}{\R{表}{ひょう}のあいだを\R{直線}{ちょく|せん}で\R{読}{よ}む：\R{内挿}{ない|そう}}

\leavevmode\label{bk:3-5:naisou}3-4では、\R{表}{ひょう}にない\R{値}{あたい}を、グラフの\R{線}{せん}から\R{読}{よ}みとりました（\R{内挿}{ない|そう}）。\R{傾}{かたむ}きを\R{使}{つか}うと、グラフをかかずに\R{計算}{けい|さん}で\R{求}{もと}められます。

\noindent\begin{minipage}[c]{0.52\linewidth}
ある\R{日}{ひ}の\R{気温}{き|おん}は、9\R{時}{じ}に\qty{18}{\degreeCelsius}、12\R{時}{じ}に\qty{24}{\degreeCelsius}でした。10\R{時}{じ}の\R{気温}{き|おん}は、\R{何度}{なん|ど}くらいだったでしょう。\par\vspace{2pt}
\R{気温}{き|おん}は\R{少}{すこ}しずつ\R{変}{か}わるので、9\R{時}{じ}から12\R{時}{じ}までは、2つの\R{点}{てん}を\R{結}{むす}ぶ\R{直線}{ちょく|せん}にそって\R{変}{か}わったと\R{考}{かんが}えます。
\end{minipage}\hfill
\begin{minipage}[c]{0.44\linewidth}\centering
\begin{tikzpicture}
\begin{axis}[bkplot, width=56mm, height=48mm, xmin=8, xmax=13.4, ymin=14, ymax=27, xtick={9,10,11,12,13}, ytick={16,18,...,26}, xlabel={\R{時}{じ}}, ylabel={\unit{\degreeCelsius}}, axis lines=left]
  \addplot[bkcurve2, domain=9:12] {18+2*(x-9)};
  \addplot[only marks, mark=*, mark size=1.8pt, bkInk] coordinates {(9,18) (12,24)};
  \draw[bkBlue, line width=0.7pt, densely dashed] (axis cs:10,14) -- (axis cs:10,20) -- (axis cs:8,20);
  \addplot[only marks, mark=o, mark size=1.8pt, bkBlue, line width=0.8pt] coordinates {(10,20)};
  \node[anchor=south east, font=\sffamily\scriptsize, text=bkBlue!80!black, inner sep=1pt] at (axis cs:9.95,20.3) {？};
\end{axis}
\end{tikzpicture}
\end{minipage}

\par\vspace{4pt}
\noindent ① \R{傾}{かたむ}き（1\R{時間}{じ|かん}あたりに\R{上}{あ}がる\R{温度}{おん|ど}）は、
\[ m=\frac{\qty{24}{\degreeCelsius}-\qty{18}{\degreeCelsius}}{\qty{12}{h}-\qty{9}{h}}=\frac{\qty{6}{\degreeCelsius}}{\qty{3}{h}}=\qty{2}{\degreeCelsius/h} \]
② 9\R{時}{じ}から10\R{時}{じ}まで、1\R{時間}{じ|かん}たったので、$\qty{2}{\degreeCelsius/h}\times\qty{1}{h}=\qty{2}{\degreeCelsius}$ \R{上}{あ}がります。$18+2=20$ で、10\R{時}{じ}は\R{約}{やく}\qty{20}{\degreeCelsius}です。

これは、\R{点}{てん} $(9,\ 18)$ と\R{傾}{かたむ}き2の\R{直線}{ちょく|せん} $y=18+2(x-9)$（\pageref{bk:3-5:tenkatamuki}ページの $y=y_1+m(x-x_1)$）に、$x=10$ を\R{入}{い}れたのと\R{同}{おな}じです。\R{表}{ひょう}のあいだの\R{値}{あたい}を、2つの\R{点}{てん}を\R{結}{むす}ぶ\R{直線}{ちょく|せん}で\R{求}{もと}めることを、\textbf{\R{直線}{ちょく|せん}\R{内挿}{ない|そう}}ともいいます。

\textbf{\R{外}{そと}へのばすのは\R{危}{あぶ}ない（\R{外挿}{がい|そう}、3-4）}\quad \R{同}{おな}じ\R{直線}{ちょく|せん}を21\R{時}{じ}までのばすと、
\[ y=18+2\times(21-9)=18+24=42 \]
で、\qty{42}{\degreeCelsius}になってしまいます。\R{夜}{よる}の\R{気温}{き|おん}がそんなに\R{高}{たか}いはずはありません。\R{気温}{き|おん}は\R{昼}{ひる}すぎから\R{下}{さ}がるので、\R{朝}{あさ}の「1\R{時間}{じ|かん}に\qty{2}{\degreeCelsius}」は、\R{夜}{よる}までは\R{続}{つづ}きません。\R{直線}{ちょく|せん}で\R{読}{よ}んでよいのは、はかった2つの\R{点}{てん}のあいだの、\R{近}{ちか}い\R{所}{ところ}だけです。

\bksubsection{9}{どこから\R{上}{うえ}になる？：1\R{次不等式}{じ|ふ|とう|しき}の\R{入}{い}り\R{口}{ぐち}}

\leavevmode\label{bk:3-5:futou}\pageref{bk:3-5:ichiji}ページの\R{水}{みず}そう $y=2x+5$ で、\R{水}{みず}が\qty{15}{L}より\R{多}{おお}くなるのは、\R{何分後}{なん|ぷん|ご}より\R{後}{あと}でしょうか。「より\R{多}{おお}い」は、\R{不等号}{ふ|とう|ごう}で $y>15$ と\R{書}{か}けます（2-2）。

\noindent\begin{minipage}[c]{0.5\linewidth}\centering
\begin{tikzpicture}
\begin{axis}[bkplot, width=60mm, height=52mm, xmin=0, xmax=8.2, ymin=0, ymax=23, xtick={1,...,8}, ytick={5,10,15,20}, xlabel={$x$（\R{分}{ふん}）}, ylabel={$y$（L）}]
  \addplot[bkRed, line width=1pt, densely dashed, domain=0:8] {15};
  \addplot[bkcurve, domain=0:5] {2*x+5};
  \addplot[bkcurve3, domain=5:8] {2*x+5};
  \addplot[only marks, mark=o, mark size=2pt, bkInk, line width=0.8pt] coordinates {(5,15)};
  \node[anchor=north, font=\sffamily\scriptsize, text=bkRed!80!black, inner sep=1pt] at (axis cs:2,14.4) {$y=15$};
  \node[anchor=east, font=\sffamily\scriptsize, text=bkGreen!50!black, inner sep=1pt] at (axis cs:6.4,19.5) {15より\R{上}{うえ}};
\end{axis}
\end{tikzpicture}
\end{minipage}\hfill
\begin{minipage}[c]{0.47\linewidth}
\textbf{グラフで}\quad $y=15$ の\R{水平}{すい|へい}な\R{線}{せん}を\R{引}{ひ}きます。\R{青}{あお}い\R{直線}{ちょく|せん}は、$x=5$ でこの\R{線}{せん}と\R{交}{まじ}わり、それより\R{右}{みぎ}では\R{線}{せん}より\R{上}{うえ}（\R{緑}{みどり}）にあります。だから、5\R{分}{ふん}より\R{後}{あと}です。\R{右上}{みぎ|あ}がりの\R{直線}{ちょく|せん}は、\R{一度}{いち|ど}こえたら、もう\R{下}{した}にはもどりません。
\end{minipage}

\par\vspace{4pt}
\textbf{\R{式}{しき}で}\quad $2x+5>15$ のように、\R{不等号}{ふ|とう|ごう}を\R{使}{つか}った\R{式}{しき}を\textbf{\R{不等式}{ふ|とう|しき}}といい、$x$ が1\R{次式}{じ|しき}なら\textbf{1\R{次不等式}{じ|ふ|とう|しき}}といいます。\bkidx{いちじふとうしき}{1\R{次不等式}{じ|ふ|とう|しき}}\R{方程式}{ほう|てい|しき}（2-16）と\R{同}{おな}じように\R{解}{と}けます。
\begin{align*}
2x+5&>15\\
2x&>15-5 && \text{\small（5を\R{移項}{い|こう}する）}\\
2x&>10\\
x&>5 && \text{\small（\R{両辺}{りょう|へん}を2でわる）}
\end{align*}
\R{不等式}{ふ|とう|しき}は、\R{片方}{かた|ほう}に\R{傾}{かたむ}いた\R{天}{てん}びんです。\R{両方}{りょう|ほう}の\R{皿}{さら}から\R{同}{おな}じ5をとっても、\R{重}{おも}いほうは\R{重}{おも}いままです。\R{両方}{りょう|ほう}を\R{半分}{はん|ぶん}にしても、\R{重}{おも}いほうは\R{重}{おも}いままです。だから、\R{不等号}{ふ|とう|ごう}の\R{向}{む}きはそのままです。

\textbf{マイナスでかける・わると、\R{向}{む}きが\R{反対}{はん|たい}になる}\quad \R{電池}{でん|ち}の $y=-20x+100$ で、\R{残}{のこ}りが30\%より\R{少}{すく}なくなるのはいつからでしょう。$-20x+100<30$ を\R{解}{と}きます。100を\R{移項}{い|こう}すると $-20x<-70$。ここで\R{両辺}{りょう|へん}を $-20$ でわります。\R{数直線}{すう|ちょく|せん}で、マイナスをかけると\R{何}{なに}が\R{起}{お}きるかを\R{見}{み}てみましょう。

\begin{bkfigure}
\begin{tikzpicture}[x=1mm, y=1mm]
  \draw[bkInk, line width=0.7pt, -{Stealth[length=4pt]}] (0,0) -- (112,0);
  \foreach \k in {-4,...,4} {\draw[bkInk, line width=0.6pt] (56+\k*12,-1.4) -- (56+\k*12,1.4);
    \node[bk tick, anchor=north] at (56+\k*12,-2.4) {$\k$};}
  \fill[bkBlue] (80,0) circle (1.6); \fill[bkBlue] (92,0) circle (1.6);
  \fill[bkOrange] (32,0) circle (1.6); \fill[bkOrange] (20,0) circle (1.6);
  \node[font=\sffamily\footnotesize, text=bkBlue!80!black, anchor=south] at (86,3) {$2<3$};
  \node[font=\sffamily\footnotesize, text=bkOrange!80!black, anchor=south] at (26,3) {$-2>-3$};
  \draw[bkGray, line width=0.6pt, -{Stealth[length=3pt]}] (79,4) to[bend right=25] node[above, font=\sffamily\scriptsize, text=bkGray, inner sep=1pt] {$\times(-1)$} (33,4);
\end{tikzpicture}
\bkfigcaption{$-1$ をかけると、0をはさんで\R{左右}{さ|ゆう}が\R{入}{い}れかわる。\R{右}{みぎ}にあった3が、\R{左}{ひだり}の $-3$ になる。}
\end{bkfigure}

$2<3$ の\R{両辺}{りょう|へん}に $-1$ をかけると、$-2$ と $-3$ になります。\R{数直線}{すう|ちょく|せん}では $-2$ のほうが\R{右}{みぎ}なので、$-2>-3$ です（2-2）。\R{大小}{だい|しょう}が\R{入}{い}れかわりました。マイナスの\R{数}{かず}でかけたり、わったりすると、\textbf{\R{不等号}{ふ|とう|ごう}の\R{向}{む}きが\R{反対}{はん|たい}}になります。
\begin{align*}
-20x+100&<30\\
-20x&<-70 && \text{\small（100を\R{移項}{い|こう}する）}\\
x&>3.5 && \text{\small（\R{両辺}{りょう|へん}を $-20$ でわる。\R{向}{む}きが\R{反対}{はん|たい}）}
\end{align*}
3.5\R{時間}{じ|かん}より\R{後}{あと}です。\R{数}{かず}で\R{確}{たし}かめると、$x=4$ なら $y=-80+100=20$ で30より\R{少}{すく}ない。$x=3$ なら $y=-60+100=40$ で30より\R{多}{おお}い。\R{右下}{みぎ|さ}がりのグラフなので、\R{右}{みぎ}へ\R{行}{い}くほど\R{低}{ひく}くなる、というグラフの\R{見方}{み|かた}とも\R{合}{あ}っています。

\begin{taisetsu}
\R{傾}{かたむ}き＝たての\R{変化量}{へん|か|りょう}÷\R{横}{よこ}の\R{変化量}{へん|か|りょう}＝$\dfrac{\Delta y}{\Delta x}=\dfrac{y_2-y_1}{x_2-x_1}$。\R{右上}{みぎ|あ}がりなら\R{正}{せい}、\R{右下}{みぎ|さ}がりなら\R{負}{ふ}、\R{水平}{すい|へい}なら0。\par\vspace{2pt}
1\R{次関数}{じ|かん|すう} $y=ax+b$：$a$ が\R{傾}{かたむ}き（\R{変化}{へん|か}の\R{割合}{わり|あい}。いつも\R{同}{おな}じ）、$b$ が\R{切片}{せっ|ぺん}（$x=0$ のときの $y$）。\par\vspace{2pt}
$x$-$t$ グラフの\R{傾}{かたむ}き＝\R{速}{はや}さ。\R{傾}{かたむ}きの\R{単位}{たん|い}は、たての\R{単位}{たん|い}÷\R{横}{よこ}の\R{単位}{たん|い}。\par\vspace{2pt}
\R{傾}{かたむ}きが $m$ で、\R{点}{てん} $(x_1,\ y_1)$ を\R{通}{とお}る\R{直線}{ちょく|せん}の\R{式}{しき}は $y-y_1=m(x-x_1)$。
\end{taisetsu}

\begin{kotoba}
\small
\bkword{\R{傾}{かたむ}き}{\R{直線}{ちょく|せん}の\R{急}{きゅう}さ。たての\R{変化量}{へん|か|りょう}÷\R{横}{よこ}の\R{変化量}{へん|か|りょう}。\R{道路}{どう|ろ}ではこう\R{配}{ばい}ともいう。}
\bkword{1\R{次関数}{じ|かん|すう}}{$y=ax+b$（$a\neq0$）の\R{形}{かたち}で\R{書}{か}ける\R{関数}{かん|すう}。グラフは\R{直線}{ちょく|せん}。}
\bkword{\R{変化}{へん|か}の\R{割合}{わり|あい}}{$x$ の\R{変化量}{へん|か|りょう}に\R{対}{たい}する $y$ の\R{変化量}{へん|か|りょう}の\R{割合}{わり|あい} $\dfrac{\Delta y}{\Delta x}$。1\R{次関数}{じ|かん|すう}ではいつも $a$。}
\bkword{\R{切片}{せっ|ぺん}}{グラフが $y$ \R{軸}{じく}と\R{交}{まじ}わる\R{点}{てん}の $y$ \R{座標}{ざ|ひょう}。$y=ax+b$ の $b$。}
\bkword{1\R{次不等式}{じ|ふ|とう|しき}}{$2x+5>15$ のような、1\R{次式}{じ|しき}の\R{不等式}{ふ|とう|しき}。マイナスでかける・わると、\R{不等号}{ふ|とう|ごう}の\R{向}{む}きが\R{反対}{はん|たい}になる。}
\end{kotoba}

%% ---------------- 例題 ----------------
\bksection{reidai}{\R{例題}{れい|だい}でたしかめよう}

\begin{reidai}[2\R{点}{てん}から\R{傾}{かたむ}きを\R{求}{もと}める]
\R{次}{つぎ}の2\R{点}{てん}を\R{通}{とお}る\R{直線}{ちょく|せん}の\R{傾}{かたむ}きを\R{求}{もと}め、\R{右上}{みぎ|あ}がりか\R{右下}{みぎ|さ}がりかを\R{答}{こた}えましょう。\par
（1）\nobreak $(1,\ 3)$、$(4,\ 9)$\qquad（2）\nobreak $(-2,\ 5)$、$(2,\ -3)$\par
\kangaekata たても\R{横}{よこ}も「2つめ－1つめ」でひきます。マイナスの\R{数}{かず}をひくときは、かっこをつけます（2-2）。\par
\shiki（1）\nobreak $\dfrac{9-3}{4-1}=\dfrac{6}{3}=2$\qquad（2）\nobreak $\dfrac{-3-5}{2-(-2)}=\dfrac{-8}{4}=-2$\quad \R{分母}{ぶん|ぼ}は $2+2=4$\par
\kotae（1）\nobreak 2、\R{右上}{みぎ|あ}がり\quad（2）\nobreak $-2$、\R{右下}{みぎ|さ}がり
\end{reidai}

\begin{reidai}[1\R{次関数}{じ|かん|すう}の\R{傾}{かたむ}き・\R{切片}{せっ|ぺん}・\R{変化量}{へん|か|りょう}]
1\R{次関数}{じ|かん|すう} $y=-3x+6$ の（1）\nobreak \R{傾}{かたむ}きと\R{切片}{せっ|ぺん}、（2）\nobreak $x$ が1から4まで\R{変}{か}わるときの $y$ の\R{変化量}{へん|か|りょう}を\R{求}{もと}めましょう。\par
\kangaekata $y=ax+b$ の $a$ が\R{傾}{かたむ}き、$b$ が\R{切片}{せっ|ぺん}。$a=-3$ はマイナスごと\R{読}{よ}みます。（2）は $y$ の\R{値}{あたい}を2つ\R{計算}{けい|さん}して、あと－まえ。\par
\shiki（1）\nobreak $a=-3$、$b=6$\par
\hspace*{2.6\zw}（2）\nobreak $x=1$ のとき\quad $y=-3\times1+6=3$\par
\hspace*{2.6\zw}\phantom{（2）}$x=4$ のとき\quad $y=-3\times4+6=-6$\par
\hspace*{2.6\zw}\phantom{（2）}$\Delta y=-6-3=-9$\par
\hspace*{2.6\zw}\R{確}{たし}かめ：\R{変化}{へん|か}の\R{割合}{わり|あい}は $-3$ なので、$\Delta y=-3\times\Delta x=-3\times3=-9$\par
\kotae（1）\nobreak \R{傾}{かたむ}き $-3$、\R{切片}{せっ|ぺん}6\quad（2）\nobreak $-9$（9へる）
\end{reidai}

\begingroup\tcbset{equal height group=bkthrfjA}%
\noindent\begin{minipage}[t]{0.49\linewidth}
\begin{jibun}\raggedright
2\R{点}{てん}を\R{通}{とお}る\R{直線}{ちょく|せん}の\R{傾}{かたむ}きを\R{求}{もと}め、\R{右上}{みぎ|あ}がりか\R{右下}{みぎ|さ}がりかを\R{答}{こた}えましょう。（1）\nobreak $(2,\ 1)$、$(5,\ 10)$\quad（2）\nobreak $(-3,\ 4)$、$(1,\ -4)$
\bkjans{（1）\nobreak $\dfrac{10-1}{5-2}=\dfrac{9}{3}=3$。\R{右上}{みぎ|あ}がり。（2）\nobreak $\dfrac{-4-4}{1-(-3)}=\dfrac{-8}{4}=-2$。\R{右下}{みぎ|さ}がり。}
\end{jibun}
\end{minipage}\hfill
\begin{minipage}[t]{0.49\linewidth}
\begin{jibun}\raggedright
1\R{次関数}{じ|かん|すう} $y=2x-4$ の（1）\nobreak \R{傾}{かたむ}きと\R{切片}{せっ|ぺん}、（2）\nobreak $x$ が1から5まで\R{変}{か}わるときの $y$ の\R{変化量}{へん|か|りょう}を\R{求}{もと}めましょう。
\bkjans{（1）\nobreak \R{傾}{かたむ}き2、\R{切片}{せっ|ぺん} $-4$。（2）\nobreak $x=1$ で $y=-2$、$x=5$ で $y=6$。$\Delta y=6-(-2)=8$（\R{確}{たし}かめ：$2\times4=8$）。}
\end{jibun}
\end{minipage}\par
\endgroup

\begin{reidai}[\R{式}{しき}からグラフをかく]
\noindent\begin{minipage}[t]{0.6\linewidth}\raggedright
\R{次}{つぎ}の1\R{次関数}{じ|かん|すう}のグラフを、ノートにかきましょう。\par
（1）\nobreak $y=\dfrac{1}{2}x-1$\qquad（2）\nobreak $y=-2x+3$\par
\kangaekata ① \R{切片}{せっ|ぺん}の\R{点}{てん} ② \R{傾}{かたむ}きの\R{分}{ぶん}だけ\R{進}{すす}んだ\R{点}{てん} ③ 2\R{点}{てん}を\R{通}{とお}る\R{直線}{ちょく|せん}、の\R{順}{じゅん}にかきます。\par
\shiki（1）\nobreak $(0,\ -1)$ から「\R{右}{みぎ}へ2・\R{上}{うえ}へ1」で $(2,\ 0)$\par
\hspace*{2.6\zw}（2）\nobreak $(0,\ 3)$ から「\R{右}{みぎ}へ1・\R{下}{した}へ2」で $(1,\ 1)$
\end{minipage}\hfill
\begin{minipage}[t]{0.37\linewidth}\vspace{0pt}\centering
\begin{tikzpicture}
\begin{axis}[bkplot, width=46mm, height=44mm, xmin=-2.4, xmax=3.6, ymin=-2.6, ymax=4.2, xtick={-2,1,2,3}, ytick={-2,1,2,3,4}, xlabel={$x$}, ylabel={$y$}]
  \addplot[bkcurve, domain=-2.2:3.4] {0.5*x-1};
  \addplot[bkcurve2, domain=-0.5:2.6] {-2*x+3};
  \addplot[only marks, mark=*, mark size=1.6pt, bkInk] coordinates {(0,-1) (2,0) (0,3) (1,1)};
  \node[anchor=south, font=\sffamily\scriptsize, text=bkBlue!80!black, inner sep=1pt] at (axis cs:3.0,0.75) {（1）};
  \node[anchor=east, font=\sffamily\scriptsize, text=bkOrange!80!black, inner sep=1pt] at (axis cs:2.0,-1.9) {（2）};
\end{axis}
\end{tikzpicture}
\end{minipage}\par
\kotae \R{図}{ず}のとおり。（1）\nobreak は $(0,\ -1)$ と $(2,\ 0)$、（2）\nobreak は $(0,\ 3)$ と $(1,\ 1)$ を\R{通}{とお}る\R{直線}{ちょく|せん}
\end{reidai}

\begin{jibun}
\R{次}{つぎ}の1\R{次関数}{じ|かん|すう}のグラフを、ノートにかきましょう。（1）\nobreak $y=\dfrac{1}{3}x+1$\quad（2）\nobreak $y=-x+2$
\bkjans{（1）\nobreak $(0,\ 1)$ と、\R{右}{みぎ}へ3・\R{上}{うえ}へ1の $(3,\ 2)$ を\R{通}{とお}る\R{直線}{ちょく|せん}。（2）\nobreak $(0,\ 2)$ と、\R{右}{みぎ}へ1・\R{下}{した}へ1の $(1,\ 1)$ を\R{通}{とお}る\R{直線}{ちょく|せん}。}
\end{jibun}

\begin{reidai}[$x$-$t$ グラフから\R{速}{はや}さを\R{求}{もと}める]
\noindent\begin{minipage}[t]{0.5\linewidth}
ソラさんが\R{家}{いえ}から\R{公園}{こう|えん}までジョギングをして、\R{少}{すこ}し\R{休}{やす}み、\R{家}{いえ}へ\R{帰}{かえ}ったときの $x$-$t$ グラフです。0〜300\R{秒}{びょう}、300〜400\R{秒}{びょう}、400〜600\R{秒}{びょう}の\R{速}{はや}さ（\R{向}{む}きもふくめて）を\R{求}{もと}めましょう。\par
\kangaekata \R{傾}{かたむ}きが\R{速}{はや}さです。\R{点}{てん}の\R{座標}{ざ|ひょう}を\R{読}{よ}み、$\dfrac{\Delta x}{\Delta t}$ を\R{単位}{たん|い}つきで\R{計算}{けい|さん}します。\R{水平}{すい|へい}なら0、\R{右下}{みぎ|さ}がりなら\R{負}{ふ}です。
\end{minipage}\hfill
\begin{minipage}[t]{0.46\linewidth}\vspace{0pt}\centering
\begin{tikzpicture}
\begin{axis}[bkplot, width=49mm, height=40mm, xmin=0, xmax=680, ymin=0, ymax=800, xtick={300,400,600}, ytick={300,600}, xlabel={$t$（s）}, ylabel={$x$（m）}, tick label style={font=\sffamily\scriptsize}]
  \addplot[bkcurve] coordinates {(0,0) (300,600) (400,600) (600,0)};
  \addplot[only marks, mark=*, mark size=1.5pt, bkInk] coordinates {(0,0) (300,600) (400,600) (600,0)};
\end{axis}
\end{tikzpicture}
\end{minipage}\par
\shiki 0〜300\R{秒}{びょう}：$\dfrac{\qty{600}{m}-\qty{0}{m}}{\qty{300}{s}-\qty{0}{s}}=\dfrac{\qty{600}{m}}{\qty{300}{s}}=\qty{2}{m/s}$\par
\hspace*{2.6\zw}300〜400\R{秒}{びょう}：$\dfrac{\qty{600}{m}-\qty{600}{m}}{\qty{400}{s}-\qty{300}{s}}=\dfrac{\qty{0}{m}}{\qty{100}{s}}=\qty{0}{m/s}$\par
\hspace*{2.6\zw}400〜600\R{秒}{びょう}：$\dfrac{\qty{0}{m}-\qty{600}{m}}{\qty{600}{s}-\qty{400}{s}}=\dfrac{\qty{-600}{m}}{\qty{200}{s}}=\qty{-3}{m/s}$\par
\kotae 0〜300\R{秒}{びょう}は\qty{2}{m/s}、300〜400\R{秒}{びょう}は\qty{0}{m/s}（\R{止}{と}まっている）、400〜600\R{秒}{びょう}は\qty{-3}{m/s}（\R{家}{いえ}のほうへ\R{秒速}{びょう|そく}\qty{3}{m}。\R{帰}{かえ}りのほうが\R{速}{はや}い）
\end{reidai}

\begin{jibun}
\noindent\begin{minipage}[t]{0.5\linewidth}
\R{右}{みぎ}の $x$-$t$ グラフは、\R{犬}{いぬ}の\R{散歩}{さん|ぽ}のようすです。0〜200\R{秒}{びょう}、200〜300\R{秒}{びょう}、300〜500\R{秒}{びょう}の\R{速}{はや}さ（\R{向}{む}きもふくめて）を\R{求}{もと}めましょう。
\end{minipage}\hfill
\begin{minipage}[t]{0.46\linewidth}\vspace{0pt}\centering
\begin{tikzpicture}
\begin{axis}[bkplot, width=49mm, height=38mm, xmin=0, xmax=560, ymin=0, ymax=1000, xtick={200,300,500}, ytick={400,800}, xlabel={$t$（s）}, ylabel={$x$（m）}]
  \addplot[bkcurve] coordinates {(0,0) (200,800) (300,800) (500,0)};
  \addplot[only marks, mark=*, mark size=1.5pt, bkInk] coordinates {(0,0) (200,800) (300,800) (500,0)};
\end{axis}
\end{tikzpicture}
\end{minipage}
\bkjans{0〜200\R{秒}{びょう}：$\dfrac{800-0}{200-0}=4$ で\qty{4}{m/s}。200〜300\R{秒}{びょう}：$\dfrac{800-800}{300-200}=0$ で\qty{0}{m/s}（\R{止}{と}まっている）。300〜500\R{秒}{びょう}：$\dfrac{0-800}{500-300}=-4$ で\qty{-4}{m/s}（\R{反対}{はん|たい}\R{向}{む}きに\qty{4}{m/s}）。}
\end{jibun}

\begin{reidai}[2\R{点}{てん}を\R{通}{とお}る\R{直線}{ちょく|せん}の\R{式}{しき}]
\R{次}{つぎ}の2\R{点}{てん}を\R{通}{とお}る\R{直線}{ちょく|せん}の\R{式}{しき}を、$y=ax+b$ の\R{形}{かたち}で\R{求}{もと}めましょう。\par
（1）\nobreak $(1,\ 5)$、$(3,\ 11)$\qquad（2）\nobreak $(-1,\ 7)$、$(2,\ 1)$\par
\kangaekata ① \R{傾}{かたむ}き $a$ を\R{求}{もと}める ② 1\R{点}{てん}を\R{代入}{だい|にゅう}して $b$ を\R{求}{もと}める ③ もう1\R{点}{てん}で\R{確}{たし}かめる。\par
\shiki（1）\nobreak $a=\dfrac{11-5}{3-1}=\dfrac{6}{2}=3$\par
\hspace*{2.6\zw}\phantom{（1）}$(1,\ 5)$ を\R{代入}{だい|にゅう}：$5=3\times1+b$、$b=2$。\R{確}{たし}かめ：$3\times3+2=11$\par
\hspace*{2.6\zw}（2）\nobreak $a=\dfrac{1-7}{2-(-1)}=\dfrac{-6}{3}=-2$\par
\hspace*{2.6\zw}\phantom{（2）}$(2,\ 1)$ を\R{代入}{だい|にゅう}：$1=-2\times2+b$、$b=5$。\R{確}{たし}かめ：$-2\times(-1)+5=7$\par
\kotae（1）\nobreak $y=3x+2$\quad（2）\nobreak $y=-2x+5$
\end{reidai}

\begin{jibun}
2\R{点}{てん}を\R{通}{とお}る\R{直線}{ちょく|せん}の\R{式}{しき}を、$y=ax+b$ の\R{形}{かたち}で\R{求}{もと}めましょう。（1）\nobreak $(2,\ 7)$、$(4,\ 11)$\quad（2）\nobreak $(-2,\ 8)$、$(1,\ 2)$
\bkjans{（1）\nobreak $a=\dfrac{11-7}{4-2}=2$。$7=2\times2+b$ より $b=3$。$y=2x+3$。（2）\nobreak $a=\dfrac{2-8}{1-(-2)}=\dfrac{-6}{3}=-2$。$2=-2\times1+b$ より $b=4$。$y=-2x+4$。}
\end{jibun}

\begin{column}{\R{電気料金}{でん|き|りょう|きん}も1\R{次関数}{じ|かん|すう}}
\R{家}{いえ}の\R{電気料金}{でん|き|りょう|きん}は、\R{使}{つか}わなくてもはらう「\R{基本料金}{き|ほん|りょう|きん}」と、\R{使}{つか}った\R{量}{りょう}に\R{比例}{ひ|れい}する\R{料金}{りょう|きん}の\R{和}{わ}です。\R{使}{つか}った\R{量}{りょう}の\R{単位}{たん|い}は kWh（キロワット\R{時}{じ}）です。\R{基本料金}{き|ほん|りょう|きん}が900\R{円}{えん}で、1\,kWh\R{使}{つか}うごとに30\R{円}{えん}かかるなら、$x$〔kWh〕\R{使}{つか}った\R{月}{つき}の\R{料金}{りょう|きん} $y$ \R{円}{えん}は $y=30x+900$ です。\R{傾}{かたむ}き30は「1\,kWhあたり30\R{円}{えん}」、\R{切片}{せっ|ぺん}900は\R{基本料金}{き|ほん|りょう|きん}です。300\,kWh\R{使}{つか}った\R{月}{つき}は $30\times300+900=9900$ で9900\R{円}{えん}。
\end{column}



\begin{reidai}[\R{表}{ひょう}のあいだの\R{値}{あたい}を\R{求}{もと}める（\R{内挿}{ない|そう}）]
やかんの\R{水}{みず}をあたためると、2\R{分後}{ふん|ご}に\qty{32}{\degreeCelsius}、6\R{分後}{ぷん|ご}に\qty{52}{\degreeCelsius}でした。このあいだは\R{直線}{ちょく|せん}で\R{変}{か}わったとして、3\R{分後}{ぷん|ご}と5\R{分後}{ふん|ご}の\R{温度}{おん|ど}を\R{求}{もと}めましょう。\par
\kangaekata ① 1\R{分}{ぷん}あたりに\R{上}{あ}がる\R{温度}{おん|ど}（\R{傾}{かたむ}き）を\R{求}{もと}める ② 2\R{分後}{ふん|ご}から\R{何分}{なん|ぷん}たったかをかけて、32にたす。\par
\shiki \R{傾}{かたむ}き：$\dfrac{\qty{52}{\degreeCelsius}-\qty{32}{\degreeCelsius}}{\qty{6}{min}-\qty{2}{min}}=\dfrac{\qty{20}{\degreeCelsius}}{\qty{4}{min}}=\qty{5}{\degreeCelsius/min}$\par
\hspace*{2.6\zw}3\R{分後}{ぷん|ご}：$32+5\times(3-2)=32+5=37$\par
\hspace*{2.6\zw}5\R{分後}{ふん|ご}：$32+5\times(5-2)=32+15=47$\par
\kotae 3\R{分後}{ぷん|ご}は\R{約}{やく}\qty{37}{\degreeCelsius}、5\R{分後}{ふん|ご}は\R{約}{やく}\qty{47}{\degreeCelsius}
\end{reidai}

\begin{reidai}[\R{場面}{ば|めん}から1\R{次関数}{じ|かん|すう}の\R{式}{しき}を\R{作}{つく}る]
ろうそくに\R{火}{ひ}をつけると、4\R{分後}{ふん|ご}に\qty{16}{cm}、10\R{分後}{ぷん|ご}に\qty{13}{cm}になりました。\R{一定}{いっ|てい}の\R{速}{はや}さで\R{短}{みじか}くなるとして、$x$ \R{分後}{ふん|ご}の\R{長}{なが}さを $y$〔cm〕とします。（1）\nobreak $y$ を $x$ の\R{式}{しき}で\R{表}{あらわ}しましょう。（2）\nobreak \R{火}{ひ}をつける\R{前}{まえ}の\R{長}{なが}さは\R{何}{なん}cmですか。（3）\nobreak \R{燃}{も}えつきるのは\R{何分後}{なん|ぷん|ご}ですか。\par
\kangaekata 「\R{一定}{いっ|てい}の\R{速}{はや}さ」なので1\R{次関数}{じ|かん|すう}で、\R{点}{てん} $(4,\ 16)$、$(10,\ 13)$ を\R{通}{とお}ります。\R{火}{ひ}をつける\R{前}{まえ}は $x=0$、\R{燃}{も}えつきるのは $y=0$ のときです。\par
\shiki（1）\nobreak $a=\dfrac{13-16}{10-4}=\dfrac{-3}{6}=-0.5$\par
\hspace*{2.6\zw}\phantom{（1）}$16=-0.5\times4+b$ より $b=16+2=18$\par
\hspace*{2.6\zw}（2）\nobreak $x=0$ を\R{入}{い}れて $y=18$\par
\hspace*{2.6\zw}（3）\nobreak $0=-0.5x+18$。\R{移項}{い|こう}して $0.5x=18$。$x=18\div0.5=36$\par
\kotae（1）\nobreak $y=-0.5x+18$\quad（2）\nobreak \qty{18}{cm}\quad（3）\nobreak 36\R{分後}{ぷん|ご}
\end{reidai}

\begingroup\tcbset{equal height group=bkthrfjC}%
\noindent\begin{minipage}[t]{0.49\linewidth}
\begin{jibun}\raggedright
ある\R{日}{ひ}の\R{気温}{き|おん}は、6\R{時}{じ}に\qty{12}{\degreeCelsius}、10\R{時}{じ}に\qty{20}{\degreeCelsius}でした。このあいだは\R{直線}{ちょく|せん}で\R{変}{か}わったとして、7\R{時}{じ}と9\R{時}{じ}の\R{気温}{き|おん}を\R{求}{もと}めましょう。
\bkjans{\R{傾}{かたむ}きは $\dfrac{20-12}{10-6}=\dfrac{8}{4}=2$ で\qty{2}{\degreeCelsius/h}。7\R{時}{じ}は $12+2\times1=14$ で\R{約}{やく}\qty{14}{\degreeCelsius}、9\R{時}{じ}は $12+2\times3=18$ で\R{約}{やく}\qty{18}{\degreeCelsius}。}
\end{jibun}
\end{minipage}\hfill
\begin{minipage}[t]{0.49\linewidth}
\begin{jibun}\raggedright
\R{水}{みず}そうの\R{水}{みず}をぬいていくと、2\R{分後}{ふん|ご}に\qty{32}{L}、5\R{分後}{ふん|ご}に\qty{20}{L}になりました。\R{一定}{いっ|てい}の\R{速}{はや}さでへるとして、$x$ \R{分後}{ふん|ご}の\R{水}{みず}の\R{量}{りょう}を $y$〔L〕とします。（1）\nobreak $y$ を $x$ の\R{式}{しき}で\R{表}{あらわ}しましょう。（2）\nobreak はじめの\R{量}{りょう}は\R{何}{なん}Lですか。（3）\nobreak \R{空}{から}になるのは\R{何分後}{なん|ぷん|ご}ですか。
\bkjans{（1）\nobreak $a=\dfrac{20-32}{5-2}=\dfrac{-12}{3}=-4$。$32=-4\times2+b$ より $b=40$。$y=-4x+40$。（2）\nobreak \qty{40}{L}。（3）\nobreak $0=-4x+40$、$4x=40$、10\R{分後}{ぷん|ご}。}
\end{jibun}
\end{minipage}\par
\endgroup

\begin{reidai}[1\R{点}{てん}と\R{傾}{かたむ}きから\R{直線}{ちょく|せん}の\R{式}{しき}]
\R{次}{つぎ}の\R{直線}{ちょく|せん}の\R{式}{しき}を、$y-y_1=m(x-x_1)$ を\R{使}{つか}って\R{作}{つく}り、$y=ax+b$ の\R{形}{かたち}に\R{直}{なお}しましょう。\par
（1）\nobreak \R{傾}{かたむ}き3で、\R{点}{てん} $(2,\ 5)$ を\R{通}{とお}る\qquad（2）\nobreak \R{傾}{かたむ}き $-\dfrac{1}{2}$ で、\R{点}{てん} $(-2,\ 4)$ を\R{通}{とお}る\par
\kangaekata $m$、$x_1$、$y_1$ に\R{数}{かず}を\R{入}{い}れます。$x_1$ がマイナスなら、かっこをつけて\R{入}{い}れます。\R{右辺}{う|へん}を\R{展開}{てん|かい}して、\R{移項}{い|こう}します。\par
\shiki\par
\noindent\begin{minipage}[t]{0.42\linewidth}
（1）\nobreak $y-5=3(x-2)$\par
\hspace*{2\zw}$y-5=3x-6$\par
\hspace*{2\zw}$y=3x-6+5$\par
\hspace*{2\zw}$y=3x-1$
\end{minipage}\hfill
\begin{minipage}[t]{0.56\linewidth}
（2）\nobreak $y-4=-\dfrac{1}{2}\{x-(-2)\}=-\dfrac{1}{2}(x+2)$\par
\hspace*{2\zw}$y-4=-\dfrac{1}{2}x-1$\par
\hspace*{2\zw}$y=-\dfrac{1}{2}x-1+4=-\dfrac{1}{2}x+3$
\end{minipage}\par
\R{確}{たし}かめ：（1）\nobreak $x=2$ で $3\times2-1=5$\par
\hspace*{5\zw}（2）\nobreak $x=-2$ で $-\dfrac{1}{2}\times(-2)+3=1+3=4$\par
\kotae（1）\nobreak $y=3x-1$\quad（2）\nobreak $y=-\dfrac{1}{2}x+3$
\end{reidai}

\begin{reidai}[1\R{次不等式}{じ|ふ|とう|しき}を\R{解}{と}く]
\R{次}{つぎ}の\R{不等式}{ふ|とう|しき}を\R{解}{と}きましょう。\quad（1）\nobreak $3x+6>21$\quad（2）\nobreak $-20x+100<30$\par
\kangaekata \R{方程式}{ほう|てい|しき}と\R{同}{おな}じように\R{移項}{い|こう}して、$x$ の\R{係数}{けい|すう}（$x$ にかかっている\R{数}{かず}）でわります。（2）は $-20$ でわるので、\R{不等号}{ふ|とう|ごう}の\R{向}{む}きが\R{反対}{はん|たい}になります。\par
\shiki（1）\nobreak $3x>21-6$、$3x>15$、$x>5$\par
\hspace*{2.6\zw}\phantom{（1）}3（プラス）でわったので、\R{向}{む}きはそのまま\par
\hspace*{2.6\zw}（2）\nobreak $-20x<30-100$、$-20x<-70$、$x>3.5$\par
\hspace*{2.6\zw}\phantom{（2）}$-20$（マイナス）でわったので、\R{向}{む}きが\R{反対}{はん|たい}\par
\hspace*{2.6\zw}\R{確}{たし}かめ：（2）\nobreak $x=4$ なら $-80+100=20$ で、30より\R{小}{ちい}さい\par
\kotae（1）\nobreak $x>5$\quad（2）\nobreak $x>3.5$
\end{reidai}

\begingroup\tcbset{equal height group=bkthrfjD}%
\noindent\begin{minipage}[t]{0.49\linewidth}
\begin{jibun}\raggedright
$y-y_1=m(x-x_1)$ を\R{使}{つか}って\R{直線}{ちょく|せん}の\R{式}{しき}を\R{作}{つく}り、$y=ax+b$ に\R{直}{なお}しましょう。（1）\nobreak \R{傾}{かたむ}き2で $(3,\ 4)$ を\R{通}{とお}る\quad（2）\nobreak \R{傾}{かたむ}き $-\dfrac{1}{3}$ で $(-3,\ 5)$ を\R{通}{とお}る
\bkjans{（1）\nobreak $y-4=2(x-3)$、$y-4=2x-6$、$y=2x-2$。（2）\nobreak $y-5=-\dfrac{1}{3}(x+3)$、$y-5=-\dfrac{1}{3}x-1$、$y=-\dfrac{1}{3}x+4$。}
\end{jibun}
\end{minipage}\hfill
\begin{minipage}[t]{0.49\linewidth}
\begin{jibun}\raggedright
\R{次}{つぎ}の\R{不等式}{ふ|とう|しき}を\R{解}{と}きましょう。\par
（1）\nobreak $4x+3>19$\par
（2）\nobreak $-5x+60<20$
\bkjans{（1）\nobreak $4x>16$、$x>4$。（2）\nobreak $-5x<-40$、\R{両辺}{りょう|へん}を $-5$ でわって\R{向}{む}きを\R{反対}{はん|たい}にし、$x>8$。}
\end{jibun}
\end{minipage}\par
\endgroup

%% ---------------- まちがえやすいところ ----------------
\Needspace{14\baselineskip}
\bksection{tsumazuki}{まちがえやすいところをチェック}\label{bk:3-5:tsuma}

\begin{tsumazuki}{たてと\R{横}{よこ}を\R{逆}{ぎゃく}にわる}
\bkng{$(1,\ 3)$、$(4,\ 9)$ の\R{傾}{かたむ}きは $\dfrac{4-1}{9-3}=\dfrac{3}{6}=0.5$}
\bkok{$\dfrac{9-3}{4-1}=\dfrac{6}{3}=2$（たての\R{変化量}{へん|か|りょう}を\R{上}{うえ}に）}
\R{傾}{かたむ}きは「\R{横}{よこ}に1\R{進}{すす}むと、たてに\R{何}{なん}\R{上}{あ}がるか」。たてを\R{横}{よこ}でわります。\R{坂}{さか}の「\R{高}{たか}さ÷\R{横}{よこ}」を\R{思}{おも}い\R{出}{だ}そう。
\end{tsumazuki}

\begin{tsumazuki}{ひく\R{順番}{じゅん|ばん}がたてと\R{横}{よこ}でちがう}
\bkng{$(1,\ 3)$、$(4,\ 9)$ の\R{傾}{かたむ}きは $\dfrac{9-3}{1-4}=\dfrac{6}{-3}=-2$}
\bkok{たても\R{横}{よこ}も「2つめ－1つめ」。$\dfrac{9-3}{4-1}=2$}
たては「2つめ－1つめ」、\R{横}{よこ}は「1つめ－2つめ」とまぜたので、\R{符号}{ふ|ごう}が\R{逆}{ぎゃく}になりました。グラフをかくと\R{右上}{みぎ|あ}がりなので、\R{傾}{かたむ}きはプラスのはずです。
\end{tsumazuki}

\begin{tsumazuki}{\R{数}{かず}だけの\R{項}{こう}を\R{傾}{かたむ}きと\R{読}{よ}む}
\bkng{$y=5-2x$ の\R{傾}{かたむ}きは5}
\bkok{$y=-2x+5$ と\R{並}{なら}べかえると、\R{傾}{かたむ}きは $-2$、\R{切片}{せっ|ぺん}は5}
\R{傾}{かたむ}きは「$x$ にかかっている\R{数}{かず}」です。\R{書}{か}いてある\R{順番}{じゅん|ばん}ではありません。$x$ が1ふえると、$y$ は2へります。
\end{tsumazuki}

\begin{tsumazuki}{マイナスの\R{座標}{ざ|ひょう}を、かっこなしで\R{入}{い}れる}
\bkng{\R{傾}{かたむ}き2で $(-1,\ 3)$ を\R{通}{とお}る\R{直線}{ちょく|せん}は $y-3=2(x-1)$}
\bkok{$y-3=2\{x-(-1)\}=2(x+1)$}
$x_1=-1$ をかっこごと\R{入}{い}れます。×の\R{式}{しき}に $x=-1$ を\R{入}{い}れると $y-3=2\times(-2)$ で $y=-1$ になり、$(-1,\ 3)$ を\R{通}{とお}りません。
\end{tsumazuki}

\begin{tsumazuki}{マイナスでわったのに、\R{不等号}{ふ|とう|ごう}をそのままにする}
\bkng{$-2x>6$ だから $x>-3$}
\bkok{$-2$ でわるので\R{向}{む}きが\R{反対}{はん|たい}。$x<-3$}
$x=0$ で\R{確}{たし}かめると、$-2\times0=0$ は6より\R{大}{おお}きくありません。でも、×の\R{答}{こた}え $x>-3$ には0が\R{入}{はい}ってしまいます。$x=-4$ なら $-2\times(-4)=8>6$ で、○の\R{答}{こた}えに\R{合}{あ}います。
\end{tsumazuki}

\begin{column}{\R{蒸気表}{じょう|き|ひょう}と\R{直線内挿}{ちょく|せん|ない|そう}}
\R{発電所}{はつ|でん|しょ}や\R{工場}{こう|じょう}のボイラーを\R{設計}{せっ|けい}する\R{人}{ひと}は、\R{水}{みず}がふっとうする\R{温度}{おん|ど}と\R{圧力}{あつ|りょく}（おす\R{強}{つよ}さ。2-10の kPa）の\R{表}{ひょう}（\R{蒸気表}{じょう|き|ひょう}）を\R{使}{つか}います。\R{表}{ひょう}には\qty{100}{\degreeCelsius}で\qty{101.42}{kPa}、\qty{110}{\degreeCelsius}で\qty{143.38}{kPa}のように、とびとびの\R{温度}{おん|ど}しかのっていません。\qty{105}{\degreeCelsius}の\R{値}{あたい}がほしいときは、\R{直線内挿}{ちょく|せん|ない|そう}をします。\R{傾}{かたむ}きは $(143.38-101.42)\div10=4.196$（1\unit{\degreeCelsius}あたりの kPa）なので、
\bkthrfline{$101.42+4.196\times5=101.42+20.98=122.40$}
\R{約}{やく}\qty{122.4}{kPa}です。\R{本当}{ほん|とう}の\R{値}{あたい}は\qty{120.90}{kPa}で、ずれは\R{約}{やく}\qty{1.5}{kPa}（\R{約}{やく}1\%）。この\R{関係}{かん|けい}のグラフは\R{少}{すこ}し\R{曲}{ま}がっているので、\R{直線}{ちょく|せん}で\R{読}{よ}むと\R{少}{すこ}しずれます。そこで、くわしい\R{蒸気表}{じょう|き|ひょう}は\R{温度}{おん|ど}を\R{細}{こま}かくきざんで、\R{内挿}{ない|そう}のずれを\R{小}{ちい}さくしています。
\end{column}



\begin{column}{スマートフォンの\R{充電}{じゅう|でん}と\R{傾}{かたむ}き}
\noindent\begin{minipage}[c]{0.58\linewidth}
\R{充電}{じゅう|でん}\R{中}{ちゅう}の\R{電池}{でん|ち}の\R{残}{のこ}り（\%）を\R{時間}{じ|かん}のグラフにすると、\R{多}{おお}くのスマートフォンでは、はじめは\R{傾}{かたむ}きが\R{大}{おお}きく（\R{速}{はや}く\R{充電}{じゅう|でん}）、80\%をこえるころから\R{傾}{かたむ}きが\R{小}{ちい}さく（ゆっくり\R{充電}{じゅう|でん}）なります。\R{満}{まん}タン\R{近}{ちか}くで\R{電流}{でん|りゅう}を\R{強}{つよ}く\R{流}{なが}しつづけると\R{電池}{でん|ち}がいたむので、\R{電子回路}{でん|し|かい|ろ}が\R{傾}{かたむ}き（\R{充電}{じゅう|でん}の\R{速}{はや}さ）を\R{調節}{ちょう|せつ}しているのです。
\end{minipage}\hfill
\begin{minipage}[c]{0.39\linewidth}\centering
\begin{tikzpicture}
\begin{axis}[bkplot, width=50mm, height=40mm, xmin=0, xmax=100, ymin=0, ymax=115, xtick={30,60,90}, ytick={50,100}, xlabel={\R{分}{ふん}}, ylabel={\%}, tick label style={font=\sffamily\scriptsize}]
  \addplot[bkcurve, domain=0:40] {2*x};
  \addplot[bkcurve, domain=40:90, samples=60] {100-20*exp(-(x-40)/10)};
\end{axis}
\end{tikzpicture}
\end{minipage}
\par\vspace{2pt}
グラフが\R{曲}{ま}がっていても、2つの\R{時刻}{じ|こく}のあいだの\R{傾}{かたむ}きは $\dfrac{\Delta y}{\Delta x}$ で\R{計算}{けい|さん}できます。それは、そのあいだの「\R{平均}{へい|きん}の\R{充電}{じゅう|でん}の\R{速}{はや}さ」です（5-1）。
\end{column}

%% ---------------- 物理・工学 ----------------
\begin{butsuri}
\par\addvspace{3pt}\textbf{\R{道路}{どう|ろ}と\R{線路}{せん|ろ}のこう\R{配}{ばい}}\par\noindent
\begin{minipage}[c]{0.76\linewidth}
\R{山}{やま}の\R{道路}{どう|ろ}には、\R{右}{みぎ}のような「10\%」の\R{標識}{ひょう|しき}があります。\R{横}{よこ}に\qty{100}{m}\R{進}{すす}むと\qty{10}{m}\R{上}{あ}がる、つまり\R{傾}{かたむ}き $\dfrac{10}{100}=0.1$ のこう\R{配}{ばい}です。\R{鉄道}{てつ|どう}は\R{車輪}{しゃ|りん}がすべりやすいので、もっとゆるやかにつくります。1000\R{分}{ぶん}の\R{何}{なん}かを‰（パーミル）で\R{表}{あらわ}し、\R{日本}{に|ほん}でいちばん\R{急}{きゅう}な\R{登山}{と|ざん}\R{電車}{でん|しゃ}でも80‰（\R{傾}{かたむ}き0.08）です。
\end{minipage}\hfill
\begin{minipage}[c]{0.2\linewidth}\centering
\begin{tikzpicture}[x=1mm, y=1mm]
  \fill[bkYellow, rounded corners=1mm, rotate around={45:(0,0)}] (-9,-9) rectangle (9,9);
  \draw[bkInk, line width=0.8pt, rounded corners=1mm, rotate around={45:(0,0)}] (-8,-8) rectangle (8,8);
  \fill[bkInk] (-6.5,-3) -- (6.5,-3) -- (6.5,3) -- cycle;
  \node[font=\sffamily\bfseries\footnotesize, text=bkInk] at (-1.5,4.5) {10\%};
\end{tikzpicture}
\end{minipage}
\par\vspace{2pt}
\R{車}{くるま}いす\R{用}{よう}のスロープは、\R{建物}{たて|もの}のバリアフリーのきまりで、$\dfrac{1}{12}$（\R{約}{やく}0.083）より\R{急}{きゅう}にしないことになっています。やってみようのスロープ（$\dfrac{30}{360}=\dfrac{1}{12}$）は、ちょうどこのきまりの\R{急}{きゅう}さです。

\par\addvspace{3pt}\textbf{\R{地図}{ち|ず}の\R{等高線}{とう|こう|せん}と\R{坂}{さか}の\R{急}{きゅう}さ}\par
\noindent\begin{minipage}[c]{0.55\linewidth}\raggedright
\R{地図}{ち|ず}では、\R{同}{おな}じ\R{高}{たか}さの\R{所}{ところ}を\R{結}{むす}んだ\R{線}{せん}（\R{等高線}{とう|こう|せん}）で、\R{山}{やま}の\R{形}{かたち}を\R{表}{あらわ}します。2\R{万}{まん}5\R{千}{せん}\R{分}{ぶん}の1の\R{地図}{ち|ず}では、\R{等高線}{とう|こう|せん}は\R{高}{たか}さ\qty{10}{m}ごとに\R{引}{ひ}かれています。\R{地図}{ち|ず}の\R{上}{うえ}で\R{線}{せん}と\R{線}{せん}の\R{間}{ま}が\qty{4}{mm}なら、\R{本当}{ほん|とう}の\R{横}{よこ}の\R{長}{なが}さは $\qty{4}{mm}\times25000=\qty{100000}{mm}=\qty{100}{m}$ なので、\R{傾}{かたむ}きは $\dfrac{10}{100}=0.1$ です。
\end{minipage}\hfill
\begin{minipage}[c]{0.43\linewidth}\centering
\begin{tikzpicture}[x=0.85mm, y=0.85mm]
  % 左がわは線がつまっていて急、右がわは線の間が広くてゆるやか
  \foreach \k/\c in {0/12, 1/22, 2/32, 3/44, 4/58} {\fill[bkGreen!\c] ({-3*\k},0) ellipse ({24-5*\k} and {15-3*\k});}
  \foreach \k in {0,...,4} {\draw[bkGreen!50!black, line width=0.5pt] ({-3*\k},0) ellipse ({24-5*\k} and {15-3*\k});}
  \draw[bkBlue, line width=0.9pt, -{Stealth[length=3pt]}] (-30,0) -- (-17,0);
  \node[font=\sffamily\scriptsize, text=bkBlue!80!black, anchor=north] at (-16,-16) {\R{左}{ひだり}から：\R{急}{きゅう}};
  \draw[bkRed, line width=0.9pt, -{Stealth[length=3pt]}] (29,0) -- (-8.5,0);
  \node[font=\sffamily\scriptsize, text=bkRed!80!black, anchor=north] at (16,-16) {\R{右}{みぎ}から：ゆるやか};
\end{tikzpicture}
\end{minipage}
\par\vspace{2pt}
\R{線}{せん}の\R{間}{ま}がせまい\R{所}{ところ}ほど、\R{同}{おな}じ\qty{10}{m}を\R{短}{みじか}い\R{横}{よこ}の\R{長}{なが}さで\R{上}{あ}がるので、\R{坂}{さか}は\R{急}{きゅう}です。\R{山}{やま}では、\R{進}{すす}む\R{向}{む}きによって\R{坂}{さか}の\R{急}{きゅう}さがちがいます。\R{東}{ひがし}へ\R{進}{すす}むときと\R{北}{きた}へ\R{進}{すす}むときの\R{傾}{かたむ}きは、8-2で\R{学}{まな}びます。
\end{butsuri}







\begin{column}{\R{温度}{おん|ど}センサーは2\R{点}{てん}で\R{合}{あ}わせる}
\R{電子}{でん|し}\R{温度計}{おん|ど|けい}の\R{中}{なか}の\R{温度}{おん|ど}センサーは、\R{温度}{おん|ど}に\R{応}{おう}じた\R{電圧}{でん|あつ}を\R{出}{だ}します。\R{電圧}{でん|あつ}は\R{温度}{おん|ど}の1\R{次関数}{じ|かん|すう}になるように\R{作}{つく}られていますが、1つ1つ\R{少}{すこ}しずつちがうので、\R{工場}{こう|じょう}では2つの\R{温度}{おん|ど}で\R{電圧}{でん|あつ}をはかって\R{直線}{ちょく|せん}を\R{決}{き}めます（\R{校正}{こう|せい}）。たとえば、\qty{0}{\degreeCelsius}で\qty{0.50}{V}、\qty{100}{\degreeCelsius}で\qty{1.50}{V}なら、\R{傾}{かたむ}きは $(1.50-0.50)\div100=0.01$ で、1\unit{\degreeCelsius}あたり\qty{0.01}{V}です。\R{温度}{おん|ど}を $x$〔\unit{\degreeCelsius}〕、\R{電圧}{でん|あつ}を $y$〔V〕とすると $y=0.01x+0.50$。\R{電圧}{でん|あつ}が\qty{0.75}{V}と\R{出}{で}たら、$0.75=0.01x+0.50$ から $0.01x=0.25$、$x=25$ で、\qty{25}{\degreeCelsius}とわかります。\R{工場}{こう|じょう}の\R{計器}{けい|き}の\R{点検}{てん|けん}でも、このように2\R{点}{てん}で\R{直線}{ちょく|せん}を\R{合}{あ}わせる\R{作業}{さ|ぎょう}をします。
\end{column}

%% ---------------- 章末問題 ----------------
\bkendproblems
\bkproblemgroup{\R{基本}{き|ほん}}{2}

\bkproblem \R{次}{つぎ}の2\R{点}{てん}を\R{通}{とお}る\R{直線}{ちょく|せん}の\R{傾}{かたむ}きを\R{求}{もと}め、\R{右上}{みぎ|あ}がりか\R{右下}{みぎ|さ}がりかを\R{答}{こた}えましょう。（1）\nobreak $(1,\ 2)$、$(4,\ 11)$\quad（2）\nobreak $(-1,\ 6)$、$(3,\ -2)$
\bkpans{（1）\nobreak $\dfrac{11-2}{4-1}=\dfrac{9}{3}=3$、\R{右上}{みぎ|あ}がり。（2）\nobreak $\dfrac{-2-6}{3-(-1)}=\dfrac{-8}{4}=-2$、\R{右下}{みぎ|さ}がり。}

\bkproblem 1\R{次関数}{じ|かん|すう} $y=-2x+7$ の（1）\nobreak \R{傾}{かたむ}きと\R{切片}{せっ|ぺん}、（2）\nobreak $x$ が2から5まで\R{変}{か}わるときの $y$ の\R{変化量}{へん|か|りょう}を\R{求}{もと}めましょう。
\bkpans{（1）\nobreak \R{傾}{かたむ}き $-2$、\R{切片}{せっ|ぺん}7。（2）\nobreak $x=2$ で $y=3$、$x=5$ で $y=-3$。$\Delta y=-3-3=-6$（$-2\times3=-6$）。}

\bkproblem \R{次}{つぎ}の1\R{次関数}{じ|かん|すう}のグラフを、ノートにかきましょう。（1）\nobreak $y=3x-2$\quad（2）\nobreak $y=-\dfrac{1}{2}x+3$
\bkpans{（1）\nobreak $(0,\ -2)$ と $(1,\ 1)$ を\R{通}{とお}る\R{直線}{ちょく|せん}。（2）\nobreak $(0,\ 3)$ と、\R{右}{みぎ}へ2・\R{下}{した}へ1の $(2,\ 2)$ を\R{通}{とお}る\R{直線}{ちょく|せん}。}

\bkproblem \noindent\begin{minipage}[t]{0.48\linewidth}
\R{右}{みぎ}の $x$-$t$ グラフは、\R{電車}{でん|しゃ}の\R{模型}{も|けい}の\R{動}{うご}きです。0〜40\R{秒}{びょう}、40〜100\R{秒}{びょう}、100〜140\R{秒}{びょう}の\R{速}{はや}さ（\R{向}{む}きもふくめて）を\R{求}{もと}めましょう。
\end{minipage}\hfill
\begin{minipage}[t]{0.44\linewidth}\vspace{0pt}\centering
\begin{tikzpicture}
\begin{axis}[bkplot, width=42mm, height=38mm, xmin=0, xmax=160, ymin=0, ymax=12, xtick={40,100,140}, ytick={5,10}, xlabel={$t$（s）}, ylabel={$x$（m）}, tick label style={font=\sffamily\scriptsize}]
  \addplot[bkcurve] coordinates {(0,0) (40,10) (100,10) (140,0)};
  \addplot[only marks, mark=*, mark size=1.4pt, bkInk] coordinates {(0,0) (40,10) (100,10) (140,0)};
\end{axis}
\end{tikzpicture}
\end{minipage}
\bkpans{0〜40\R{秒}{びょう}：$\dfrac{10}{40}=0.25$ で\qty{0.25}{m/s}。40〜100\R{秒}{びょう}：\qty{0}{m/s}（\R{止}{と}まっている）。100〜140\R{秒}{びょう}：$\dfrac{0-10}{140-100}=-0.25$ で\qty{-0.25}{m/s}（\R{反対}{はん|たい}\R{向}{む}きに\qty{0.25}{m/s}）。}

\bkproblem \R{次}{つぎ}の2\R{点}{てん}を\R{通}{とお}る\R{直線}{ちょく|せん}の\R{式}{しき}を、$y=ax+b$ の\R{形}{かたち}で\R{求}{もと}めましょう。（1）\nobreak $(3,\ 7)$、$(5,\ 11)$\quad（2）\nobreak $(-2,\ 9)$、$(2,\ 1)$
\bkpans{（1）\nobreak $a=\dfrac{11-7}{5-3}=2$、$7=2\times3+b$ より $b=1$。$y=2x+1$。（2）\nobreak $a=\dfrac{1-9}{2-(-2)}=-2$、$1=-2\times2+b$ より $b=5$。$y=-2x+5$。}

\bkproblem \R{次}{つぎ}の\R{直線}{ちょく|せん}の\R{式}{しき}を、$y-y_1=m(x-x_1)$ を\R{使}{つか}って\R{作}{つく}り、$y=ax+b$ の\R{形}{かたち}に\R{直}{なお}しましょう。（1）\nobreak \R{傾}{かたむ}き4で $(1,\ 3)$ を\R{通}{とお}る\quad（2）\nobreak \R{傾}{かたむ}き $-\dfrac{1}{2}$ で $(-4,\ 5)$ を\R{通}{とお}る
\bkpans{（1）\nobreak $y-3=4(x-1)$、$y=4x-1$。（2）\nobreak $y-5=-\dfrac{1}{2}(x+4)$、$y-5=-\dfrac{1}{2}x-2$、$y=-\dfrac{1}{2}x+3$。}

\bkproblem \R{水}{みず}の\R{温度}{おん|ど}が、1\R{分後}{ぷん|ご}に\qty{24}{\degreeCelsius}、5\R{分後}{ふん|ご}に\qty{40}{\degreeCelsius}でした。このあいだは\R{直線}{ちょく|せん}で\R{変}{か}わったとして、3\R{分後}{ぷん|ご}の\R{温度}{おん|ど}を\R{求}{もと}めましょう。
\bkpans{\R{傾}{かたむ}きは $\dfrac{40-24}{5-1}=4$ で\qty{4}{\degreeCelsius/min}。$24+4\times(3-1)=32$ で\R{約}{やく}\qty{32}{\degreeCelsius}。}

\bkproblem \R{次}{つぎ}の\R{不等式}{ふ|とう|しき}を\R{解}{と}きましょう。（1）\nobreak $5x-4>11$\quad（2）\nobreak $-3x+9<0$
\bkpans{（1）\nobreak $5x>15$、$x>3$。（2）\nobreak $-3x<-9$、$-3$ でわって\R{向}{む}きを\R{反対}{はん|たい}にし、$x>3$。}

\bkproblemgroup{チャレンジ}{3}

\bkproblem ばねにおもりをつるすと、3\R{個}{こ}で\qty{26}{cm}、7\R{個}{こ}で\qty{34}{cm}になりました（ばね\R{全体}{ぜん|たい}の\R{長}{なが}さ）。おもりの\R{数}{かず}を $x$ \R{個}{こ}、\R{長}{なが}さを $y$〔cm〕とします。（1）\nobreak $y$ を $x$ の\R{式}{しき}で\R{表}{あらわ}しましょう。（2）\nobreak おもりをつるす\R{前}{まえ}の\R{長}{なが}さは\R{何}{なん}cmですか。（3）\nobreak \R{長}{なが}さが\qty{40}{cm}になるのは、おもり\R{何個}{なん|こ}のときですか。
\bkpans{（1）\nobreak $a=\dfrac{34-26}{7-3}=2$、$26=2\times3+b$ より $b=20$。$y=2x+20$。（2）\nobreak \qty{20}{cm}。（3）\nobreak $40=2x+20$、$2x=20$、10\R{個}{こ}。}

\bkproblem スマートフォンの\R{電池}{でん|ち}が80\%あり、\R{動画}{どう|が}を\R{見}{み}つづけると1\R{時間}{じ|かん}に15\%ずつへります。$x$ \R{時間後}{じ|かん|ご}の\R{残}{のこ}りを $y$\,\% とします。（1）\nobreak $y$ を $x$ の\R{式}{しき}で\R{表}{あらわ}しましょう。（2）\nobreak \R{残}{のこ}りが20\%より\R{少}{すく}なくなるのは、\R{何時間後}{なん|じ|かん|ご}より\R{後}{あと}ですか。\R{不等式}{ふ|とう|しき}を\R{作}{つく}って\R{解}{と}きましょう。
\bkpans{（1）\nobreak $y=-15x+80$。（2）\nobreak $-15x+80<20$、$-15x<-60$、$-15$ でわって\R{向}{む}きを\R{反対}{はん|たい}にし、$x>4$。4\R{時間後}{じ|かん|ご}より\R{後}{あと}。}

\bkproblem \noindent\begin{minipage}[t]{0.48\linewidth}
\R{右}{みぎ}の $x$-$t$ グラフは、ミナさんが\R{家}{いえ}からポストまで\R{歩}{ある}いて\R{手紙}{て|がみ}を\R{出}{だ}し、\R{家}{いえ}へもどったときのようすです。（1）\nobreak 0〜100\R{秒}{びょう}、100〜250\R{秒}{びょう}、250〜400\R{秒}{びょう}の\R{速}{はや}さ（\R{向}{む}きもふくめて）を\R{求}{もと}めましょう。（2）\nobreak \R{行}{い}きと\R{帰}{かえ}りでは、どちらが\R{速}{はや}いですか。
\end{minipage}\hfill
\begin{minipage}[t]{0.44\linewidth}\vspace{0pt}\centering
\begin{tikzpicture}
\begin{axis}[bkplot, width=46mm, height=38mm, xmin=0, xmax=460, ymin=0, ymax=380, xtick={100,250,400}, ytick={150,300}, xlabel={$t$（s）}, ylabel={$x$（m）}, tick label style={font=\sffamily\scriptsize}]
  \addplot[bkcurve] coordinates {(0,0) (100,300) (250,300) (400,0)};
  \addplot[only marks, mark=*, mark size=1.4pt, bkInk] coordinates {(0,0) (100,300) (250,300) (400,0)};
\end{axis}
\end{tikzpicture}
\end{minipage}
\bkpans{（1）\nobreak 0〜100\R{秒}{びょう}：$\dfrac{300-0}{100-0}=3$ で\qty{3}{m/s}。100〜250\R{秒}{びょう}：\qty{0}{m/s}（\R{止}{と}まっている）。250〜400\R{秒}{びょう}：$\dfrac{0-300}{400-250}=-2$ で\qty{-2}{m/s}。（2）\nobreak \R{行}{い}き（\qty{3}{m/s}）のほうが\R{速}{はや}い。\R{帰}{かえ}りは\R{反対}{はん|たい}\R{向}{む}きに\qty{2}{m/s}。}

\bkproblem \R{傾}{かたむ}き $-3$ で、\R{点}{てん} $(-2,\ 1)$ を\R{通}{とお}る\R{直線}{ちょく|せん}の\R{式}{しき}を、$y-y_1=m(x-x_1)$ を\R{使}{つか}って\R{作}{つく}り、$y=ax+b$ の\R{形}{かたち}に\R{直}{なお}しましょう。また、$x=-2$ を\R{入}{い}れて\R{確}{たし}かめましょう。
\bkpans{$y-1=-3\{x-(-2)\}=-3(x+2)$、$y-1=-3x-6$、$y=-3x-5$。\R{確}{たし}かめ：$-3\times(-2)-5=6-5=1$。}

\begin{furikaeri}
\begin{checklist}
\item \R{傾}{かたむ}きを「たての\R{変化量}{へん|か|りょう}÷\R{横}{よこ}の\R{変化量}{へん|か|りょう}」で\R{求}{もと}め、\R{正}{せい}・\R{負}{ふ}・0がわかる。\bkfurref{bk:3-5:katamuki}
\item $y=ax+b$ の\R{傾}{かたむ}きと\R{切片}{せっ|ぺん}を\R{読}{よ}み、$x$-$t$ グラフの\R{傾}{かたむ}きを\R{速}{はや}さとして\R{使}{つか}える。\bkfurref{bk:3-5:ichiji}
\item 1\R{点}{てん}と\R{傾}{かたむ}きから、$y-y_1=m(x-x_1)$ で\R{直線}{ちょく|せん}の\R{式}{しき}を\R{作}{つく}れる。\bkfurref{bk:3-5:tenkatamuki}
\end{checklist}
\end{furikaeri}

\bkfukidashi{delta}{\R{急}{きゅう}さを\R{数}{かず}にしたら、\R{坂}{さか}も\R{速}{はや}さも\R{電池}{でん|ち}も、\R{同}{おな}じ「\R{傾}{かたむ}き」で\R{話}{はな}せたね。\R{次}{つぎ}の\R{章}{しょう}では、2\R{本}{ほん}の\R{直線}{ちょく|せん}が\R{交}{まじ}わる\R{点}{てん}を\R{求}{もと}めるよ。}
